% concatenated from arxiv_final_copy.tex
\documentclass{article}
\usepackage{arxiv}
\usepackage[utf8]{inputenc} % allow utf-8 input
\usepackage[T1]{fontenc}    % use 8-bit T1 fonts
\usepackage{hyperref, enumitem, url, amsmath, amssymb, amsthm, booktabs, float, amsfonts, subcaption, microtype, graphicx, natbib, doi, algorithm, algpseudocode}
\numberwithin{equation}{section}
%
%% ====== Macros & Notation ======
\DeclareMathOperator*{\argmin}{arg\,min}
\newcommand{\norm}[1]{\left\|#1\right\|}
\DeclareMathOperator{\dist}{dist}
\DeclareMathOperator{\Var}{Var}
\DeclareMathOperator{\Cov}{Cov}
\graphicspath{{results_noharm_piil/figures/}{results_noharm_map/figures/}}

\theoremstyle{definition}
\newtheorem{assumption}{Assumption}
\newtheorem{definition}{Definition}
\newtheorem{proposition}{Proposition}
\newtheorem{remark}{Remark}
\newtheorem{lemma}{Lemma}
\newtheorem{corollary}{Corollary}
\newtheorem{theorem}{Theorem}
\newtheorem{example}{Example}


\title{No-Harm Physics-Informed Inverse Learning with Residual-Calibrated Uncertainty}

%\date{September 9, 1985}	% Here you can change the date presented in the paper title
%\date{} 					% Or removing it

\author{\href{https://orcid.org/0000-0002-8545-1833}{\includegraphics[scale=0.06]{orcid.pdf}}
Ronald Katende \\
Department of Mathematics\\
Kabale University\\
Kikungiri Hill, Katuna Road, Kabale, Uganda \\
\texttt{rkatende@kab.ac.ug}}

\date{}


\renewcommand{\headeright}{}
\renewcommand{\undertitle}{}
\renewcommand{\shorttitle}{No-Harm Physics-Informed Inverse Learning with Residual-Calibrated Uncertainty}

\hypersetup{
pdftitle={No-Harm Physics-Informed Inverse Learning with Residual-Calibrated Uncertainty},
pdfauthor={Ronald Katende},
pdfkeywords={physics-informed learning; inverse problems; residual certification; uncertainty quantification; a posteriori error estimation; stochastic collocation; tomography; geophysical inversion; numerical methods},
}

\begin{document}
\maketitle

\begin{abstract}
Physics-informed learning is increasingly used for partial differential equation (PDE)-governed inverse problems, but its reliability remains difficult to certify. This paper develops a no-harm certification-and-selection framework for physics-informed inverse learning. A learned reconstruction is accepted only when its residual-calibrated radius is no worse than the baseline radius, namely when $$R_{\mathrm{learn}}\le R_{\mathrm{base}}+\varepsilon_{\mathrm{safe}};$$otherwise, the method returns the baseline. The certificate combines data, physics, boundary or initial-condition, and optimization residuals. Under a conditional stability estimate, these residuals yield an a posteriori reconstruction-error bound and a deterministic uncertainty radius. A high-probability certificate is also derived for physics residuals estimated from independent random collocation points. Numerical tests on Poisson source recovery, inverse heat reconstruction, limited-angle tomography, elliptic coefficient identification, and stochastic residual validation show that the selector accepts certified improvements, rejects shifted, hallucinated, or unfinished candidates, and becomes conservative in strongly ill-posed regimes. The framework is therefore a certification-and-selection layer, not another reconstruction architecture.
\end{abstract}

\keywords{physics-informed learning \and inverse problems \and residual certification \and uncertainty quantification \and a posteriori error estimation \and stochastic collocation \and tomography \and geophysical inversion \and numerical methods}

\section{Introduction}
\label{sec:introduction}

Inverse problems arise when an unknown quantity must be recovered from indirect, incomplete, or noisy observations. They occur in source recovery for partial differential equations, limited-angle tomography, medical and scientific imaging, geophysical coefficient identification, and parameter estimation in physical systems \cite{EnglHankeNeubauer1996,Kirsch2011,KaipioSomersalo2005,Hansen2010}. Their central difficulty is ill-posedness: the inverse map may fail to exist, fail to be unique, or be unstable under perturbations of the data \cite{TikhonovArsenin1977,EnglHankeNeubauer1996}. Regularization and numerical control are therefore not secondary implementation details. They are the mechanisms by which inverse reconstructions become stable enough to be interpreted.

Data-driven methods have changed the practical landscape of inverse problems. Learned reconstruction maps, learned regularizers, unrolled iterative schemes, neural operators, and plug-and-play methods can use training data to improve reconstructions in familiar regimes \cite{ArridgeMaassOktemSchonlieb2019,AdlerOktem2018,LiKovachkiAzizzadenesheli2021}. At the same time, these methods raise reliability questions that are especially serious in ill-posed settings. A reconstruction may be visually plausible or accurate on average while remaining unstable under small perturbations, changes in sampling geometry, noise shift, model mismatch, or out-of-distribution structure \cite{AntunRennaPoonAdcockHansen2020}. Artificial Intelligence (AI)-assisted inverse solvers therefore need more than accuracy metrics. They need computable reliability information.

Physics-informed learning offers a natural bridge between learned reconstruction and numerical modelling. Physics-informed neural networks (PINNs) and related methods train neural representations using observed data together with residuals of governing equations, boundary or initial conditions, and sometimes conservation or constitutive constraints \cite{RaissiPerdikarisKarniadakis2019,KarniadakisKevrekidisLuPerdikarisWangYang2021}. In inverse settings, the neural model may approximate both the state \(u\) and an unknown parameter, source, coefficient, or field \(q\). The loss typically combines a data-misfit term with a PDE residual, boundary or initial-condition residuals, and regularization or prior terms. This structure is useful because it embeds the forward physics into the learning problem instead of treating the inverse map as a purely black-box regression task.

The remaining difficulty is certification. A small training loss or a plausible reconstructed field does not by itself guarantee that the recovered parameter is reliable. Failure can occur because the data are sparse, the inverse problem is weakly identifiable, the residual is evaluated on too few collocation points, the optimizer has not reached a meaningful stationary point, or the assumed forward model does not match the data-generating system. These issues are familiar in numerical inverse problems, where residuals, stability estimates, and a posteriori error indicators are used to connect computation with reliability \cite{EnglHankeNeubauer1996,Kirsch2011,Verfuerth2013}. They are also central to uncertainty quantification, where the goal is not only to compute a reconstruction but to describe how strongly the available data support it \cite{KaipioSomersalo2005,Sullivan2015,Stuart2010}.

This paper develops a no-harm certification-and-selection framework for physics-informed inverse learning. The central principle is simple: a learned inverse reconstruction should be allowed to improve a baseline, but it should not be allowed to silently degrade reliability. Given a learned reconstruction \((\widehat u,\widehat q)\) and a baseline reconstruction \((u_{\mathrm{base}},q_{\mathrm{base}})\), the method computes a residual-calibrated reliability radius from four quantities: the data residual, the physics residual, the boundary or initial-condition residual, and an optimization residual. Under a conditional stability estimate for the inverse problem, this radius gives an a posteriori bound on the reconstruction error. The learned reconstruction is accepted only when its certified radius is no worse than the baseline radius, up to a prescribed tolerance. Otherwise, the method returns the baseline reconstruction.

The framework is deliberately conditional. It does not claim that every physics-informed neural solver is reliable, nor that residuals alone solve ill-posedness. Instead, it identifies the assumptions under which residual information becomes meaningful. If the inverse problem admits a conditional stability estimate on the admissible class and the residuals are evaluated with sufficient accuracy, then the certificate converts data misfit, physics violation, boundary violation, and optimization error into a computable reliability statement. The solver output is therefore not only a reconstruction. It is a reconstruction accompanied by a no-harm decision and an uncertainty radius.

The contributions are as follows.
\begin{enumerate}[label=(\roman*)]
\item We define a residual certificate for physics-informed inverse learning that combines data, PDE, boundary or initial-condition, and optimization residuals.
\item We prove that, under a conditional stability estimate, the certificate yields an a posteriori error bound for the recovered parameter, source, or coefficient.
\item We introduce a residual-calibrated uncertainty radius that can be reported together with a learned inverse reconstruction.
\item We propose a no-harm certification-and-selection rule that accepts a learned reconstruction only when its certified radius improves on, or matches, a baseline certificate up to a prescribed tolerance; otherwise, the framework falls back to the baseline.
\item We derive high-probability residual bounds for the case where the physics residual is estimated from random collocation points.
\item We include an optimization-residual certificate to account for approximate neural minimizers and unfinished training.
\item We design numerical tests in inverse source recovery, inverse heat reconstruction, limited-angle tomography, and geophysical coefficient identification to assess reconstruction quality, certificate sharpness, and behaviour under noise or distribution shift.
\end{enumerate}


\section{Problem setting}
\label{sec:problem-setting}

Let \(\Omega\subset\mathbb R^d\) be a bounded Lipschitz domain. Let \(\mathcal X\), \(\mathcal Q\), \(\mathcal Y\), \(\mathcal Z\), and \(\mathcal B_0\) denote Hilbert or Banach spaces appropriate for the state, the unknown parameter, the observations, the interior residual, and the boundary or initial residual.

The spaces are kept abstract because the certificate is intended to apply to several inverse-problem classes. The state norm on \(\mathcal X\) measures the physical field reconstructed by the forward model. The parameter norm on \(\mathcal Q\) is the norm in which the unknown source, coefficient, initial condition, or material field is certified. The data norm on \(\mathcal Y\) measures agreement with observations. The residual norm on \(\mathcal Z\) measures violation of the governing equation, while \(\mathcal B_0\) measures boundary or initial-condition mismatch. These choices separate the physical reconstruction error from the residual quantities used to certify it.

We consider inverse problems governed by
\begin{equation}
\mathcal N(u,q)=0
\quad \text{in } \Omega,
\label{eq:pde}
\end{equation}
together with a boundary or initial condition
\begin{equation}
\mathcal B(u,q)=0
\quad \text{on } \partial\Omega
\quad \text{or on the prescribed initial set}.
\label{eq:boundary}
\end{equation}
The observation model is
\begin{equation}
y^\delta=\mathcal H u^\dagger+e,
\qquad
\|e\|_{\mathcal Y}\le \delta,
\label{eq:observation}
\end{equation}
where \(u^\dagger\) is the exact state, \(q^\dagger\) is the exact unknown, \(\mathcal H:\mathcal X\to\mathcal Y\) is the observation operator, and \(\delta\ge0\) is a noise level. This formulation covers many standard inverse problems in which \(q\) may be a source term, coefficient field, initial condition, boundary input, conductivity, permeability, material parameter, or forcing term \cite{EnglHankeNeubauer1996,Kirsch2011,KaipioSomersalo2005}.

The objective is to recover \(q^\dagger\) from \(y^\delta\). The certificate is stated for \(q\), rather than primarily for \(u\), because inverse problems usually seek the hidden source, coefficient, initial condition, boundary input, or material parameter that generated the observed state. The state \(u\) remains essential because it links the unknown \(q\) to the observations and to the governing equation, but the reliability question is whether the recovered unknown \(q\) is trustworthy. In ill-posed settings, this recovery cannot be justified by data fit alone. A candidate pair \((u,q)\) must also be assessed through the governing equation, boundary or initial conditions, and the stability structure of the inverse map \cite{TikhonovArsenin1977,EnglHankeNeubauer1996,Hansen2010}. This is the basis for the residual certificates developed below.

\subsection{Physics-informed inverse reconstruction}

A physics-informed inverse solver represents the state and unknown quantity by trainable functions
\begin{equation}
u_\theta:\Omega\to\mathbb R^{d_u},
\qquad
q_\theta:\Omega\to\mathbb R^{d_q},
\label{eq:neural-state-parameter}
\end{equation}
or by equivalent neural, finite-dimensional, or hybrid parameterizations. Here \(d_u\) and \(d_q\) denote the pointwise dimensions of the state and unknown fields. The symbol \(m\) is reserved for the number of observations in the numerical experiments. The parameters \(\theta\) are trained by minimizing a loss of the form
\begin{equation}
\mathcal L(\theta)
=
\mathcal L_{\mathrm{data}}(\theta)
+
\alpha_{\mathrm{pde}}\mathcal L_{\mathrm{pde}}(\theta)
+
\alpha_{\mathrm{bc}}\mathcal L_{\mathrm{bc}}(\theta)
+
\alpha_{\mathrm{reg}}\mathcal L_{\mathrm{reg}}(\theta).
\label{eq:pinn-loss}
\end{equation}
The data term measures agreement with \(y^\delta\), the PDE term penalizes violation of \eqref{eq:pde}, the boundary or initial term penalizes violation of \eqref{eq:boundary}, and the regularization term encodes additional prior or numerical control.

The regularization term is used to stabilize the reconstruction procedure, but it is not automatically treated as a certificate residual. The certificate below measures quantities that can be directly checked after reconstruction: data mismatch, physics violation, boundary or initial-condition violation, and optimization stationarity. If a prior or regularization condition is intended to be part of the reliability claim, it can be added as an additional certified residual term. In this paper, regularization supports the solver, while the reported no-harm decision is based on residual-calibrated reliability.

This structure is standard in physics-informed learning and its inverse-problem variants, where the forward physics is imposed through residual terms rather than learned only from input-output samples \cite{RaissiPerdikarisKarniadakis2019,KarniadakisKevrekidisLuPerdikarisWangYang2021}.

After training, the solver returns
\begin{equation}
\widehat u=u_{\widehat\theta},
\qquad
\widehat q=q_{\widehat\theta}.
\label{eq:learned-output}
\end{equation}
The pair \((\widehat u,\widehat q)\) is a computed reconstruction, not automatically a certified one. The distinction matters because small empirical loss may still coexist with sparse data, weak identifiability, poor collocation coverage, optimizer error, or model mismatch \cite{AntunRennaPoonAdcockHansen2020,CuomoDiColaGiampaoloRozzaRaissiPiccialli2022}.

\subsection{Baseline reconstruction}

Let \((u_{\mathrm{base}},q_{\mathrm{base}})\) denote a baseline reconstruction. The baseline may be a classical variational inverse solution, a finite-element or finite-difference PDE-constrained optimization solution, a conservative regularized reconstruction, or a standard physics-informed reconstruction without the additional learned component being tested. Such baselines are common in inverse problems and PDE-constrained optimization because they provide transparent reference points for stability, regularization, and numerical comparison \cite{EnglHankeNeubauer1996,HinzePinnauUlbrichUlbrich2009}.

The baseline is not assumed to be exact or optimal. Its role is to provide a reference certificate. A learned reconstruction is accepted only if it can match or improve this reference under the residual-calibrated reliability measure introduced below. This is the no-harm principle used in this paper.


\section{Residual certificates}
\label{sec:residual-certificates}

\subsection{Data, physics, and boundary residuals}

For a candidate pair \((u,q)\), define the data residual
\begin{equation}
r_{\mathrm{data}}(u)
=
\|\mathcal H u-y^\delta\|_{\mathcal Y}.
\label{eq:data-residual}
\end{equation}
This measures agreement with the observed data. It is necessary but not sufficient for reliability, since ill-posed inverse problems may admit reconstructions that fit noisy data while remaining unstable or physically inconsistent \cite{EnglHankeNeubauer1996,Hansen2010}.

Define the physics residual by
\begin{equation}
r_{\mathrm{pde}}(u,q)
=
\|\mathcal N(u,q)\|_{\mathcal Z},
\label{eq:pde-residual}
\end{equation}
where \(\mathcal Z\) is the residual space associated with the governing equation. For strong-form physics-informed training, this may be an \(L^2(\Omega)\)-type norm estimated at collocation points. For weak or variational formulations, it may be a dual norm over a test space, as in residual-based numerical analysis and a posteriori error estimation \cite{Verfuerth2013,BangerthRannacher2003}.

Define the boundary or initial residual by
\begin{equation}
r_{\mathrm{bc}}(u,q)
=
\|\mathcal B(u,q)\|_{\mathcal B_0}.
\label{eq:boundary-residual}
\end{equation}
This term is kept separate from the interior physics residual because boundary and initial errors may have different scaling, different sampling procedures, and different effects on stability.

\subsection{Optimization residual}

Physics-informed inverse learning is usually solved by nonconvex numerical optimization. A small reported loss does not by itself show that the optimizer has reached a meaningful stationary point. We therefore include an optimization residual. For a differentiable training loss, define
\begin{equation}
r_{\mathrm{opt}}(\theta)
=
\|\nabla_\theta \mathcal L(\theta)\|.
\label{eq:optimization-residual}
\end{equation}
For constrained, projected, or nonsmooth formulations, \(r_{\mathrm{opt}}\) can be replaced by a projected-gradient norm, a Karush--Kuhn--Tucker (KKT) residual, or a subdifferential distance. Such stationarity measures are standard diagnostics in numerical optimization and are necessary when the reconstruction is produced by an iterative solver rather than by exact minimization \cite{NocedalWright2006,Bertsekas1999}.

\subsection{Total residual certificate}

The total residual certificate is defined by
\begin{equation}
\eta(u,q,\theta)
=
r_{\mathrm{data}}(u)
+
\alpha_{\mathrm{pde}} r_{\mathrm{pde}}(u,q)
+
\alpha_{\mathrm{bc}} r_{\mathrm{bc}}(u,q)
+
\alpha_{\mathrm{opt}} r_{\mathrm{opt}}(\theta),
\label{eq:total-certificate}
\end{equation}
where the weights \(\alpha_{\mathrm{pde}},\alpha_{\mathrm{bc}},\alpha_{\mathrm{opt}}\ge0\) place the residual components on comparable scales. This scaling is necessary because data residuals, PDE residuals, boundary or initial-condition residuals, and optimization residuals may have different units, dimensions, and numerical magnitudes. The weights may be chosen from dimensional scaling, residual normalization, validation data, or conservative user-specified tolerances, but they must be fixed before comparing the learned reconstruction with the baseline. The same weighting rule is used for both candidates. The role of \(\eta\) is diagnostic and certifying: it records how well the candidate reconstruction fits the data, satisfies the physics, respects boundary or initial constraints, and reflects a sufficiently solved optimization problem.

\begin{definition}[Residual-certified reconstruction]
\label{def:residual-certified}
A physics-informed inverse reconstruction \((\widehat u,\widehat q)\), produced by parameters \(\widehat\theta\), is residual-certified at level \(\rho\ge0\) if
\begin{equation}
\eta(\widehat u,\widehat q,\widehat\theta)\le \rho.
\label{eq:residual-certified}
\end{equation}
\end{definition}

This definition does not by itself imply small reconstruction error. The connection between residual size and reconstruction error requires a stability property of the inverse problem. That property is stated next.

\section{Conditional stability and a posteriori error control}
\label{sec:conditional-stability}

\subsection{Conditional inverse stability}

The link between residual size and reconstruction error requires a stability property of the inverse problem. Such stability cannot be assumed globally in most inverse problems. Many inverse maps are compact, smoothing, or only conditionally stable on restricted admissible classes. Depending on the problem, stability may be Lipschitz, H\"older, logarithmic, or valid only under source, geometric, monotonicity, or regularity constraints \cite{EnglHankeNeubauer1996,Isakov2006,Yamamoto2009,Kirsch2011}. The certification results below therefore do not assume that residuals alone solve ill-posedness. They assume an explicit conditional stability estimate and then show how computable residuals produce an a posteriori reliability bound.

\begin{assumption}[Conditional stability]
\label{ass:conditional-stability}
There exist a constant \(C_{\mathrm{stab}}>0\), an exponent \(p\in(0,1]\), and an admissible set
\[
\mathcal A\subset \mathcal X\times\mathcal Q
\]
such that, for every \((u,q)\in\mathcal A\),
\begin{equation}
\|q-q^\dagger\|_{\mathcal Q}
\le
C_{\mathrm{stab}}
\left(
\|\mathcal H u-\mathcal H u^\dagger\|_{\mathcal Y}
+
\|\mathcal N(u,q)\|_{\mathcal Z}
+
\|\mathcal B(u,q)\|_{\mathcal B_0}
\right)^p .
\label{eq:conditional-stability}
\end{equation}
\end{assumption}

Assumption~\ref{ass:conditional-stability} is deliberately problem-dependent. The admissible set \(\mathcal A\) may encode regularity, bounds, positivity, structural constraints, parameter ranges, or a model class in which the inverse problem is stable. This is standard in inverse-problem analysis, where useful stability estimates are usually tied to an admissible class rather than to the whole ambient space \cite{Isakov2006,EnglHankeNeubauer1996}. In the present paper, the stability estimate is used as the bridge between numerical residuals and certified reconstruction error.

\subsection{A posteriori reconstruction bound}

\begin{theorem}[Residual a posteriori error bound]
\label{thm:residual-error-bound}
Assume the observation model \eqref{eq:observation} and Assumption~\ref{ass:conditional-stability}. Let \((\widehat u,\widehat q)\in\mathcal A\) be any computed physics-informed reconstruction. Then
\begin{equation}
\|\widehat q-q^\dagger\|_{\mathcal Q}
\le
C_{\mathrm{stab}}
\left(
r_{\mathrm{data}}(\widehat u)
+
r_{\mathrm{pde}}(\widehat u,\widehat q)
+
r_{\mathrm{bc}}(\widehat u,\widehat q)
+
\delta
\right)^p .
\label{eq:aposteriori-error-bound}
\end{equation}
\end{theorem}

\begin{proof}
By the triangle inequality and the noise model \eqref{eq:observation},
\begin{equation}
\|\mathcal H\widehat u-\mathcal H u^\dagger\|_{\mathcal Y}
\le
\|\mathcal H\widehat u-y^\delta\|_{\mathcal Y}
+
\|y^\delta-\mathcal H u^\dagger\|_{\mathcal Y}
\le
r_{\mathrm{data}}(\widehat u)+\delta .
\label{eq:data-triangle-bound}
\end{equation}
Since \((\widehat u,\widehat q)\in\mathcal A\), Assumption~\ref{ass:conditional-stability} applies. Substituting \eqref{eq:data-triangle-bound} and the definitions of \(r_{\mathrm{pde}}\) and \(r_{\mathrm{bc}}\) into \eqref{eq:conditional-stability} gives \eqref{eq:aposteriori-error-bound}.
\end{proof}

Theorem~\ref{thm:residual-error-bound} is an a posteriori statement. It does not depend on how \((\widehat u,\widehat q)\) was obtained. The pair may come from a physics-informed neural network, a hybrid neural-finite-element method, a classical PDE-constrained optimizer, or another numerical inverse solver. What matters for the bound is admissibility, data consistency, physics consistency, boundary consistency, and the noise level.

\begin{definition}[Residual-calibrated uncertainty radius]
\label{def:uncertainty-radius}
For any admissible candidate \((u,q)\in\mathcal A\), define
\begin{equation}
R_\delta(u,q)
=
C_{\mathrm{stab}}
\left(
r_{\mathrm{data}}(u)
+
r_{\mathrm{pde}}(u,q)
+
r_{\mathrm{bc}}(u,q)
+
\delta
\right)^p .
\label{eq:uncertainty-radius}
\end{equation}
The associated certified uncertainty set is
\begin{equation}
\mathcal U_\delta(u,q)
=
\left\{
\widetilde q\in\mathcal Q:
\|\widetilde q-q\|_{\mathcal Q}
\le
R_\delta(u,q)
\right\}.
\label{eq:uncertainty-set}
\end{equation}
\end{definition}

For numerical certification and no-harm decisions, we also use the operational radius
\begin{equation}
R_{\delta}^{\mathrm{op}}(u,q,\theta)
=
C_{\mathrm{stab}}
\left(
r_{\mathrm{data}}(u)
+
\alpha_{\mathrm{pde}} r_{\mathrm{pde}}(u,q)
+
\alpha_{\mathrm{bc}} r_{\mathrm{bc}}(u,q)
+
\delta
+
\alpha_{\mathrm{opt}} r_{\mathrm{opt}}(\theta)
\right)^p .
\label{eq:operational-radius}
\end{equation}
The weights in \eqref{eq:operational-radius} are fixed before evaluation and reported in the validation protocol. The deterministic containment result is attached to \(R_\delta\), because \(R_\delta\) follows directly from the conditional stability estimate in Assumption~\ref{ass:conditional-stability}. The numerical no-harm decisions use \(R_{\delta}^{\mathrm{op}}\) when scaled residuals and optimization residuals are included. Thus \(R_\delta\) is the theoretical radius, while \(R_{\delta}^{\mathrm{op}}\) is the implemented decision radius. Both are based on the same no-harm logic, but the operational radius includes practical scaling and stationarity diagnostics needed in numerical learning.

The set \(\mathcal U_\delta(u,q)\) is deterministic. It is not a Bayesian credible set and should not be interpreted as a posterior distribution. It is a residual-calibrated uncertainty set obtained from an inverse stability estimate and a noise bound. Bayesian inverse-problem uncertainty and deterministic residual certification answer different questions, and they rely on different assumptions \cite{KaipioSomersalo2005,Stuart2010,Sullivan2015}.

\begin{corollary}[Certified containment]
\label{cor:certified-containment}
Under the assumptions of Theorem~\ref{thm:residual-error-bound},
\begin{equation}
q^\dagger\in\mathcal U_\delta(\widehat u,\widehat q).
\label{eq:certified-containment}
\end{equation}
\end{corollary}

\begin{proof}
Theorem~\ref{thm:residual-error-bound} gives
\[
\|\widehat q-q^\dagger\|_{\mathcal Q}
\le
R_\delta(\widehat u,\widehat q).
\]
This is exactly the membership condition \(q^\dagger\in\mathcal U_\delta(\widehat u,\widehat q)\).
\end{proof}


\section{A no-harm principle for physics-informed inverse learning}
\label{sec:no-harm}

The residual radius developed in Section~\ref{sec:conditional-stability} turns a computed inverse reconstruction into a certifiable object. We now use that radius to state the central principle of the paper. The principle is not tied to a particular neural architecture, optimizer, discretization, or inverse-problem class. It applies to any physics-informed inverse solver that returns a candidate reconstruction together with computable residuals.

\subsection{The principle}

A learned physics-informed reconstruction should not be accepted merely because it has a small training loss, a plausible visual appearance, or a lower error in a synthetic experiment where the ground truth is known. In real inverse problems, the true reconstruction error is unavailable. The decision must therefore be made from quantities available at reconstruction time: data residual, physics residual, boundary or initial-condition residual, optimization residual, noise level, and a stability estimate.

This leads to the no-harm principle used in this paper.

\begin{definition}[No-harm principle for physics-informed inverse learning]
\label{def:no-harm-principle}
Let \((u_{\mathrm{base}},q_{\mathrm{base}})\) be a baseline reconstruction and let \((\widehat u,\widehat q)\) be a learned physics-informed reconstruction. Let
\[
R_{\mathrm{base}}
\quad\text{and}\quad
R_{\mathrm{learn}}
\]
be their corresponding certified radii, computed using the same certification rule. For a prescribed tolerance \(\varepsilon_{\mathrm{safe}}\ge0\), the learned reconstruction is selected only if
\begin{equation}
R_{\mathrm{learn}}
\le
R_{\mathrm{base}}+\varepsilon_{\mathrm{safe}}.
\label{eq:no-harm-principle-rule}
\end{equation}
If \eqref{eq:no-harm-principle-rule} fails, the method falls back to the baseline reconstruction.
\end{definition}

The principle is deliberately conservative. It does not repair or post-process an unreliable learned reconstruction. It decides whether the learned output is certified enough to replace a safer reference solution. The framework is therefore a certification-and-selection layer: its output is a selected reconstruction together with the certificate and residual components that justify the selection.

\subsection{General certificate-dominance formulation}

Let \(\mathcal C\) be a class of admissible reconstruction procedures. A procedure \(A\in\mathcal C\) takes the noisy data \(y^\delta\) and returns a candidate pair
\begin{equation}
A(y^\delta)=(u_A,q_A).
\label{eq:procedure-output}
\end{equation}
Assume that each admissible candidate can be assigned a computable certificate
\begin{equation}
\mathfrak R(A)
=
\mathfrak R(u_A,q_A).
\label{eq:abstract-certificate}
\end{equation}
The certificate \(\mathfrak R\) may be the deterministic residual radius \(R_\delta\), the stochastic radius \(R_{\delta,\zeta}^{\mathrm{stoch}}\), or the operational radius \(R_\delta^{\mathrm{op}}\) used in numerical implementation. The only requirement is that the same certificate rule is used for the baseline and for the learned candidate being tested.

\begin{definition}[Certificate-valid reconstruction]
\label{def:certificate-valid-reconstruction}
A candidate reconstruction \(A(y^\delta)=(u_A,q_A)\) is certificate-valid if
\begin{equation}
\|q_A-q^\dagger\|_{\mathcal Q}
\le
\mathfrak R(A).
\label{eq:certificate-valid}
\end{equation}
\end{definition}

For the deterministic certificate, validity follows from Theorem~\ref{thm:residual-error-bound}. For the stochastic certificate, validity holds with probability at least \(1-\zeta\) under Corollary~\ref{cor:stochastic-containment}. For the operational certificate used in the numerical experiments, \(\mathfrak R\) is the reported residual-calibrated reliability radius used by the no-harm decision.

Let \(A_{\mathrm{base}}\) denote the baseline procedure and \(A_{\mathrm{learn}}\) denote the learned physics-informed procedure. Write
\begin{equation}
A_{\mathrm{base}}(y^\delta)
=
(u_{\mathrm{base}},q_{\mathrm{base}}),
\qquad
A_{\mathrm{learn}}(y^\delta)
=
(u_{\mathrm{learn}},q_{\mathrm{learn}}).
\label{eq:base-learn-procedure-outputs}
\end{equation}
Here \((u_{\mathrm{learn}},q_{\mathrm{learn}})\) is the learned output denoted elsewhere by \((\widehat u,\widehat q)\).

\begin{definition}[No-harm selector]
\label{def:no-harm-selector}
The no-harm selector \(S_{\mathrm{NH}}\) is defined by
\begin{equation}
S_{\mathrm{NH}}(y^\delta)
=
\begin{cases}
	A_{\mathrm{learn}}(y^\delta),
	&
	\text{if }
	\mathfrak R(A_{\mathrm{learn}})
	\le
	\mathfrak R(A_{\mathrm{base}})+\varepsilon_{\mathrm{safe}},
	\\[0.4em]
	A_{\mathrm{base}}(y^\delta),
	&
	\text{otherwise}.
\end{cases}
\label{eq:no-harm-selector}
\end{equation}
\end{definition}

The selected output is denoted by
\begin{equation}
(u_{\mathrm{safe}},q_{\mathrm{safe}})
=
S_{\mathrm{NH}}(y^\delta).
\label{eq:selected-safe-output}
\end{equation}

\begin{theorem}[No-harm certificate dominance]
\label{thm:no-harm-principle}
Assume that \(A_{\mathrm{base}}(y^\delta)\) and \(A_{\mathrm{learn}}(y^\delta)\) are certificate-valid in the sense of Definition~\ref{def:certificate-valid-reconstruction}. Then the no-harm selected output satisfies
\begin{equation}
\|q_{\mathrm{safe}}-q^\dagger\|_{\mathcal Q}
\le
\mathfrak R(A_{\mathrm{base}})
+
\varepsilon_{\mathrm{safe}}.
\label{eq:no-harm-dominance}
\end{equation}
If the learned reconstruction is rejected, then
\begin{equation}
\|q_{\mathrm{safe}}-q^\dagger\|_{\mathcal Q}
\le
\mathfrak R(A_{\mathrm{base}}).
\label{eq:no-harm-rejected-dominance}
\end{equation}
\end{theorem}

\begin{proof}
There are two cases.

If the learned reconstruction is selected, then
\[
q_{\mathrm{safe}}=q_{\mathrm{learn}}.
\]
By certificate validity of \(A_{\mathrm{learn}}\),
\[
\|q_{\mathrm{learn}}-q^\dagger\|_{\mathcal Q}
\le
\mathfrak R(A_{\mathrm{learn}}).
\]
Since the learned reconstruction is selected, Definition~\ref{def:no-harm-selector} gives
\[
\mathfrak R(A_{\mathrm{learn}})
\le
\mathfrak R(A_{\mathrm{base}})
+
\varepsilon_{\mathrm{safe}}.
\]
Combining the last two inequalities yields
\[
\|q_{\mathrm{safe}}-q^\dagger\|_{\mathcal Q}
=
\|q_{\mathrm{learn}}-q^\dagger\|_{\mathcal Q}
\le
\mathfrak R(A_{\mathrm{base}})
+
\varepsilon_{\mathrm{safe}}.
\]

If the learned reconstruction is rejected, then
\[
q_{\mathrm{safe}}=q_{\mathrm{base}}.
\]
By certificate validity of \(A_{\mathrm{base}}\),
\[
\|q_{\mathrm{base}}-q^\dagger\|_{\mathcal Q}
\le
\mathfrak R(A_{\mathrm{base}}).
\]
Therefore
\[
\|q_{\mathrm{safe}}-q^\dagger\|_{\mathcal Q}
=
\|q_{\mathrm{base}}-q^\dagger\|_{\mathcal Q}
\le
\mathfrak R(A_{\mathrm{base}}).
\]
This proves \eqref{eq:no-harm-rejected-dominance}. Since
\[
\mathfrak R(A_{\mathrm{base}})
\le
\mathfrak R(A_{\mathrm{base}})
+
\varepsilon_{\mathrm{safe}},
\]
the bound \eqref{eq:no-harm-dominance} also holds in the rejected case.
\end{proof}

Theorem~\ref{thm:no-harm-principle} is the formal no-harm guarantee. It does not assert that the learned reconstruction is always more accurate than the baseline. It states a certificate-dominance rule: the learned reconstruction is selected only when its available certificate is no worse than the baseline certificate, up to the prescribed tolerance. Otherwise, the method falls back to the baseline. This prevents an uncertified learned output from silently replacing a safer reference solution.

\begin{corollary}[Strict no-harm rule]
\label{cor:strict-no-harm}
If \(\varepsilon_{\mathrm{safe}}=0\), then the selected reconstruction satisfies
\begin{equation}
\|q_{\mathrm{safe}}-q^\dagger\|_{\mathcal Q}
\le
\mathfrak R(A_{\mathrm{base}}).
\label{eq:strict-no-harm}
\end{equation}
Moreover, whenever the learned reconstruction is selected,
\begin{equation}
\mathfrak R(A_{\mathrm{learn}})
\le
\mathfrak R(A_{\mathrm{base}}).
\label{eq:selected-certificate-dominance}
\end{equation}
\end{corollary}

\begin{proof}[Proof of Corollary~\ref{cor:strict-no-harm}]
Set \(\varepsilon_{\mathrm{safe}}=0\) in Theorem~\ref{thm:no-harm-principle}. Then
\[
\|q_{\mathrm{safe}}-q^\dagger\|_{\mathcal Q}
\le
\mathfrak R(A_{\mathrm{base}}),
\]
which proves \eqref{eq:strict-no-harm}.

If the learned reconstruction is selected, Definition~\ref{def:no-harm-selector} gives
\[
\mathfrak R(A_{\mathrm{learn}})
\le
\mathfrak R(A_{\mathrm{base}})
+
\varepsilon_{\mathrm{safe}}.
\]
With \(\varepsilon_{\mathrm{safe}}=0\), this becomes
\[
\mathfrak R(A_{\mathrm{learn}})
\le
\mathfrak R(A_{\mathrm{base}}),
\]
which proves \eqref{eq:selected-certificate-dominance}.
\end{proof}

\subsection{Algorithmic use}

The no-harm principle can be implemented as a wrapper around any baseline inverse solver and any learned physics-informed inverse solver. The wrapper does not modify the learned reconstruction. It evaluates the certificate, compares it with the baseline certificate, and selects the safer certified output.

\begin{algorithm}[!ht]
\caption{No-harm certification and selection for physics-informed inverse learning}
\label{alg:no-harm-piil}
\begin{algorithmic}[1]
\Require Noisy data \(y^\delta\), operators \(\mathcal N\), \(\mathcal B\), and \(\mathcal H\), baseline solver \(A_{\mathrm{base}}\), learned physics-informed solver \(A_{\mathrm{learn}}\), noise level \(\delta\), safety tolerance \(\varepsilon_{\mathrm{safe}}\), and certificate rule \(\mathfrak R\).
\State Compute the baseline reconstruction:
\[
(u_{\mathrm{base}},q_{\mathrm{base}})
\gets
A_{\mathrm{base}}(y^\delta).
\]
\State Compute the baseline residual components:
\[
r_{\mathrm{data}}(u_{\mathrm{base}}),\quad
r_{\mathrm{pde}}(u_{\mathrm{base}},q_{\mathrm{base}}),\quad
r_{\mathrm{bc}}(u_{\mathrm{base}},q_{\mathrm{base}}).
\]
\State Compute the baseline certificate:
\[
R_{\mathrm{base}}
\gets
\mathfrak R(A_{\mathrm{base}}).
\]
\State Compute the learned reconstruction:
\[
(\widehat u,\widehat q)
\gets
A_{\mathrm{learn}}(y^\delta).
\]
\State Compute the learned residual components:
\[
r_{\mathrm{data}}(\widehat u),\quad
r_{\mathrm{pde}}(\widehat u,\widehat q),\quad
r_{\mathrm{bc}}(\widehat u,\widehat q),\quad
r_{\mathrm{opt}}(\widehat\theta).
\]
\State Compute the learned certificate:
\[
R_{\mathrm{learn}}
\gets
\mathfrak R(A_{\mathrm{learn}}).
\]
\If{\(R_{\mathrm{learn}}\le R_{\mathrm{base}}+\varepsilon_{\mathrm{safe}}\)}
\State Select the learned reconstruction:
\[
(u_{\mathrm{safe}},q_{\mathrm{safe}})
\gets
(\widehat u,\widehat q).
\]
\State Set the decision label to \(\mathrm{accept}\).
\Else
\State Fall back to the baseline reconstruction:
\[
(u_{\mathrm{safe}},q_{\mathrm{safe}})
\gets
(u_{\mathrm{base}},q_{\mathrm{base}}).
\]
\State Set the decision label to \(\mathrm{reject}\).
\EndIf
\State \Return selected reconstruction \((u_{\mathrm{safe}},q_{\mathrm{safe}})\), decision label, \(R_{\mathrm{base}}\), \(R_{\mathrm{learn}}\), residual components, and selected method.
\end{algorithmic}
\end{algorithm}

In deterministic certification, \(\mathfrak R\) may be chosen as \(R_\delta\). When the physics residual is estimated from independent validation collocation points, \(\mathfrak R\) may be chosen as \(R_{\delta,\zeta}^{\mathrm{stoch}}\). In numerical implementations with scaled residuals and optimization residuals, \(\mathfrak R\) may be chosen as the operational radius \(R_\delta^{\mathrm{op}}\), provided the same operational certificate is used for both the baseline and the learned candidate.

The step-by-step application is simple. First, compute a baseline. Second, compute a learned physics-informed candidate. Third, evaluate both certificates using the same rule. Fourth, accept the learned candidate only if it passes the no-harm inequality. Fifth, report the selected reconstruction together with the residual components and the decision. The certificate is not hidden inside the solver; it is part of the reported output.

\subsection{Interpretation}

The no-harm principle separates reconstruction from certification. A learned inverse solver may produce a visually plausible reconstruction, but visual plausibility is not a reliability certificate. It may produce a lower error in a synthetic experiment, but the true error is unavailable in real applications. It may produce a small physics-informed training loss, but this can still coexist with weak identifiability, sparse observations, collocation undersampling, optimizer error, or model mismatch. The no-harm principle therefore asks a different question:
\[
\text{Is the learned output certified enough to replace the baseline?}
\]
If the answer is yes, the learned reconstruction is selected. If the answer is no, the baseline is returned.

This makes the principle architecture-independent. It can wrap physics-informed neural networks, neural operators, unrolled inverse solvers, hybrid finite-element neural methods, classical PDE-constrained optimizers, or any future inverse solver that can report the residuals needed for certification. The principle changes the role of physics-informed inverse learning from
\[
\text{reconstruct and hope}
\]
to
\[
\text{reconstruct, certify, and select}.
\]

\subsection{Connection with the numerical validation}

The numerical validation is designed to test this principle directly. In the Poisson source and geophysical coefficient experiments, the good learned candidates are selected because their certified radii are smaller than the baseline radii. Shifted, hallucinated, and unfinished learned candidates are rejected because their certificates are weaker. In the inverse heat experiment, the certificate becomes conservative as the inverse problem becomes more ill-posed. In limited-angle tomography, a learned candidate can have lower true error in hindsight but still be rejected because its certificate is worse than the baseline certificate. This is not a failure of the rule. It is the intended behaviour: the selector uses only information available at reconstruction time.

\subsection{Three demonstrative applications of the no-harm principle}
\label{subsec:no-harm-demonstrative-applications}

The no-harm principle can be used as a simple decision layer around any inverse solver that reports a certificate. This subsection gives three concrete applications. The purpose is not to introduce new experiments, but to show how the principle is applied in practice.

\subsubsection{Application 1: accepting a learned reconstruction when its certificate improves on the baseline}

Consider a baseline reconstruction \((u_{\mathrm{base}},q_{\mathrm{base}})\) and a learned reconstruction \((\widehat u,\widehat q)\). Suppose both are evaluated using the same operational certificate rule. Let
\[
R_{\mathrm{base}}=5.375,
\qquad
R_{\mathrm{learn}}=4.930,
\qquad
\varepsilon_{\mathrm{safe}}=0 .
\]
The no-harm decision condition is
\[
R_{\mathrm{learn}}
\le
R_{\mathrm{base}}+\varepsilon_{\mathrm{safe}} .
\]
Substituting the values gives
\[
4.930 \le 5.375 .
\]
The learned reconstruction is therefore selected:
\[
(u_{\mathrm{safe}},q_{\mathrm{safe}})
=
(\widehat u,\widehat q).
\]
If the certificates are valid, Theorem~\ref{thm:no-harm-principle} gives
\[
\|q_{\mathrm{safe}}-q^\dagger\|_{\mathcal Q}
=
\|\widehat q-q^\dagger\|_{\mathcal Q}
\le
R_{\mathrm{learn}}
\le
R_{\mathrm{base}} .
\]
Thus the learned reconstruction is not selected because it is produced by a neural or physics-informed solver. It is selected because its certified radius is smaller than the certified radius of the baseline. This is the intended behaviour when the learned inverse solver improves the certified reliability of the reconstruction.

\subsubsection{Application 2: rejecting a learned reconstruction even when it appears accurate}

A learned reconstruction may look plausible, or may even have lower true error in a synthetic test, while still being less reliable according to the available certificate. Suppose
\[
R_{\mathrm{base}}=6.520,
\qquad
R_{\mathrm{learn}}=7.922,
\qquad
\varepsilon_{\mathrm{safe}}=0 .
\]
Then
\[
7.922 > 6.520 ,
\]
so the no-harm condition fails. The selected output is
\[
(u_{\mathrm{safe}},q_{\mathrm{safe}})
=
(u_{\mathrm{base}},q_{\mathrm{base}}).
\]
This is not a failure of the method. The no-harm selector is not an oracle for the unknown reconstruction error. In real inverse problems, the true error
\[
\|\widehat q-q^\dagger\|_{\mathcal Q}
\]
is unavailable. The decision must therefore use the quantities available at reconstruction time. If the learned reconstruction has a larger residual-calibrated radius than the baseline, the method falls back to the baseline.

This case is important for limited-angle tomography and other ill-posed imaging problems. A learned prior may produce visually convincing structure, but visual plausibility is not a certificate. The no-harm rule requires the learned output to be supported by the same residual-calibrated reliability test used for the baseline.

\subsubsection{Application 3: using independent validation collocation points}

In physics-informed inverse learning, the physics residual is often estimated from collocation points. A training residual alone may be optimistic, especially if the same points are used for training and certification. The no-harm principle can therefore be used with an independent validation residual.

Let \(X_1,\ldots,X_M\) be independent validation collocation points. For a candidate \((u,q)\), compute
\[
\widehat r_{\mathrm{pde}}^{\,2}(u,q)
=
\frac{1}{M}
\sum_{j=1}^{M}
\left|
\mathcal N(u,q)(X_j)
\right|^2 .
\]
Under the bounded residual-envelope assumption in Section~\ref{sec:stochastic-residuals}, define
\[
\widehat r_{\mathrm{pde},\zeta}(u,q)
=
\left[
\widehat r_{\mathrm{pde}}^{\,2}(u,q)
+
B_{\mathrm{pde}}
\sqrt{
\frac{\log(1/\zeta)}{2M}
}
\right]^{1/2}.
\]
The corresponding stochastic certified radius is
\[
R_{\delta,\zeta}^{\mathrm{stoch}}(u,q)
=
C_{\mathrm{stab}}
\left(
r_{\mathrm{data}}(u)
+
\widehat r_{\mathrm{pde},\zeta}(u,q)
+
r_{\mathrm{bc}}(u,q)
+
\delta
\right)^p .
\]
This stochastic radius can be used as the certificate rule \(\mathfrak R\) in the no-harm selector. The learned reconstruction is selected only if
\[
R_{\delta,\zeta}^{\mathrm{stoch}}(\widehat u,\widehat q)
\le
R_{\delta,\zeta}^{\mathrm{stoch}}(u_{\mathrm{base}},q_{\mathrm{base}})
+
\varepsilon_{\mathrm{safe}} .
\]
Otherwise, the method returns the baseline. This gives a direct, reproducible workflow for physics-informed inverse learning: train the learned model, evaluate its residual on independent validation collocation points, compute the stochastic certificate, compare it with the baseline certificate, and select only if the no-harm inequality holds.

\subsubsection{Practical workflow}

The three cases above reduce the no-harm principle to a reproducible five-step procedure:
\begin{enumerate}[label=(\roman*)]
\item compute a baseline reconstruction and its certificate;
\item compute a learned physics-informed reconstruction and its certificate using the same rule;
\item compare the two certified radii using the no-harm inequality;
\item select the learned reconstruction only if its certified radius is no worse than the baseline radius up to \(\varepsilon_{\mathrm{safe}}\);
\item report the selected reconstruction, the two radii, the residual components, and the accept-or-fallback decision.
\end{enumerate}
The method therefore changes the operational use of physics-informed inverse learning from ``reconstruct and hope'' to ``reconstruct, certify, and select.'' 


\section{Stochastic residual estimation}
\label{sec:stochastic-residuals}

Physics-informed methods often approximate residual norms by sampling collocation points. This is computationally convenient, but it introduces statistical error into the certificate. A residual evaluated on a small collocation set may underestimate the true residual over the domain. This section converts an empirical physics residual into a high-probability upper bound and thereby defines a stochastic certificate that can be used as the rule \(\mathfrak R\) in the no-harm selector of Section~\ref{sec:no-harm}. The certificate is stated for a candidate reconstruction that is fixed before the validation collocation points are drawn. This is the correct a posteriori use case: after training, the reconstruction is frozen and then checked on an independent residual-validation set. The result is not a uniform generalization bound over all possible neural-network parameters. The result is stated for an independent validation collocation set, because using the same points for training and certification may lead to optimistic residual estimates. This separation between training error and validation error is standard in statistical learning and is important for reliable certification \cite{Vapnik1998,ShalevShwartzBenDavid2014}.

\subsection{Empirical physics residual}

Let \(X_1,\dots,X_M\) be independent collocation points sampled from a probability measure \(\nu\) on \(\Omega\). For a fixed admissible pair \((u,q)\), define the empirical squared physics residual by
\begin{equation}
\widehat r_{\mathrm{pde}}^{\,2}(u,q)
=
\frac{1}{M}
\sum_{j=1}^{M}
\left|
\mathcal N(u,q)(X_j)
\right|^2 .
\label{eq:empirical-pde-residual}
\end{equation}
The corresponding population squared physics residual is
\begin{equation}
r_{\mathrm{pde},\nu}^{2}(u,q)
=
\int_{\Omega}
\left|
\mathcal N(u,q)(x)
\right|^2
\,d\nu(x).
\label{eq:population-pde-residual}
\end{equation}
When \(\nu\) is normalized Lebesgue measure, \(r_{\mathrm{pde},\nu}\) is the normalized \(L^2\)-type physics residual. More generally, \(\nu\) may be chosen to reflect the validation distribution used for residual checking.

\begin{assumption}[Bounded residual envelope]
\label{ass:bounded-residual}
For the fixed admissible pair \((u,q)\) being certified, there exists \(B_{\mathrm{pde}}<\infty\) such that
\begin{equation}
0
\le
\left|
\mathcal N(u,q)(x)
\right|^2
\le
B_{\mathrm{pde}}
\quad
\text{for \(\nu\)-almost every } x\in\Omega .
\label{eq:bounded-residual}
\end{equation}
\end{assumption}

The bounded-envelope assumption is used only to obtain a simple concentration bound. Other concentration tools may be substituted when the residual has sub-Gaussian, sub-exponential, or problem-specific tail behaviour.

\begin{theorem}[High-probability physics-residual certificate]
\label{thm:high-prob-residual}
Let \((u,q)\) be fixed before the validation collocation points are drawn, and suppose Assumption~\ref{ass:bounded-residual} holds. Then, for any \(\zeta\in(0,1)\), with probability at least \(1-\zeta\) over the draw of \(X_1,\dots,X_M\),
\begin{equation}
r_{\mathrm{pde},\nu}^{2}(u,q)
\le
\widehat r_{\mathrm{pde}}^{\,2}(u,q)
+
B_{\mathrm{pde}}
\sqrt{
	\frac{\log(1/\zeta)}{2M}
}.
\label{eq:high-prob-residual}
\end{equation}
\end{theorem}

\begin{proof}
Define
\[
Z_j
=
\left|
\mathcal N(u,q)(X_j)
\right|^2,
\qquad
j=1,\dots,M.
\]
Since \(X_1,\dots,X_M\) are independent, the random variables \(Z_1,\dots,Z_M\) are independent. By Assumption~\ref{ass:bounded-residual},
\[
0\le Z_j\le B_{\mathrm{pde}}
\quad
\text{for each }j.
\]
Moreover,
\[
\mathbb E[Z_j]
=
\int_{\Omega}
\left|
\mathcal N(u,q)(x)
\right|^2
\,d\nu(x)
=
r_{\mathrm{pde},\nu}^{2}(u,q).
\]
Hoeffding's inequality gives, for every \(t>0\),
\[
\mathbb P\left(
r_{\mathrm{pde},\nu}^{2}(u,q)
-
\frac{1}{M}\sum_{j=1}^{M}Z_j
\ge t
\right)
\le
\exp\left(
-\frac{2Mt^2}{B_{\mathrm{pde}}^2}
\right).
\]
Choose
\[
t
=
B_{\mathrm{pde}}
\sqrt{
\frac{\log(1/\zeta)}{2M}
}.
\]
Then
\[
\exp\left(
-\frac{2Mt^2}{B_{\mathrm{pde}}^2}
\right)
=
\zeta.
\]
Hence, with probability at least \(1-\zeta\),
\[
r_{\mathrm{pde},\nu}^{2}(u,q)
-
\frac{1}{M}\sum_{j=1}^{M}Z_j
\le
B_{\mathrm{pde}}
\sqrt{
\frac{\log(1/\zeta)}{2M}
}.
\]
Using the definition of \(\widehat r_{\mathrm{pde}}^{\,2}(u,q)\) in \eqref{eq:empirical-pde-residual} gives \eqref{eq:high-prob-residual}.
\end{proof}

Define the high-probability physics residual by
\begin{equation}
\widehat r_{\mathrm{pde},\zeta}(u,q)
=
\left[
\widehat r_{\mathrm{pde}}^{\,2}(u,q)
+
B_{\mathrm{pde}}
\sqrt{
\frac{\log(1/\zeta)}{2M}
}
\right]^{1/2}.
\label{eq:hp-pde-residual}
\end{equation}
Then Theorem~\ref{thm:high-prob-residual} states that
\begin{equation}
r_{\mathrm{pde},\nu}(u,q)
\le
\widehat r_{\mathrm{pde},\zeta}(u,q)
\label{eq:hp-pde-residual-bound}
\end{equation}
with probability at least \(1-\zeta\).

\begin{remark}
Theorem~\ref{thm:high-prob-residual} certifies a fixed candidate reconstruction on an independent validation collocation set. A uniform certificate over a whole neural-network class would require capacity control, covering arguments, Rademacher complexity, PAC-Bayes bounds, or another uniform-convergence tool \cite{Vapnik1998,ShalevShwartzBenDavid2014}. The present paper uses the fixed-candidate version because it matches the a posteriori workflow: after training, the computed reconstruction is checked on an independent residual-validation set.
\end{remark}

\subsection{High-probability uncertainty radius}

Replacing the population physics residual by the high-probability upper bound gives the stochastic residual-calibrated certificate
\begin{equation}
R_{\delta,\zeta}^{\mathrm{stoch}}(u,q)
=
C_{\mathrm{stab}}
\left(
r_{\mathrm{data}}(u)
+
\widehat r_{\mathrm{pde},\zeta}(u,q)
+
r_{\mathrm{bc}}(u,q)
+
\delta
\right)^p .
\label{eq:stochastic-radius}
\end{equation}
This is the stochastic analogue of the deterministic radius \(R_\delta\). In numerical implementations, one may instead use an operational stochastic radius that also includes scaled residual components and an optimization residual, provided the same operational rule is applied to both the baseline and the learned candidate.

\begin{corollary}[High-probability certified containment]
\label{cor:stochastic-containment}
Assume that the conditional stability estimate in Assumption~\ref{ass:conditional-stability} holds with \(r_{\mathrm{pde},\nu}\) as the physics residual. Let \((\widehat u,\widehat q)\in\mathcal A\) be fixed before drawing the independent validation collocation set. Under Assumption~\ref{ass:bounded-residual}, with probability at least \(1-\zeta\),
\begin{equation}
\|\widehat q-q^\dagger\|_{\mathcal Q}
\le
R_{\delta,\zeta}^{\mathrm{stoch}}(\widehat u,\widehat q).
\label{eq:stochastic-containment}
\end{equation}
\end{corollary}

\begin{proof}
By Theorem~\ref{thm:high-prob-residual},
\[
r_{\mathrm{pde},\nu}(\widehat u,\widehat q)
\le
\widehat r_{\mathrm{pde},\zeta}(\widehat u,\widehat q)
\]
with probability at least \(1-\zeta\). On this event, the residual a posteriori bound in Theorem~\ref{thm:residual-error-bound}, with \(r_{\mathrm{pde},\nu}\) as the physics residual, gives
\[
\|\widehat q-q^\dagger\|_{\mathcal Q}
\le
C_{\mathrm{stab}}
\left(
r_{\mathrm{data}}(\widehat u)
+
r_{\mathrm{pde},\nu}(\widehat u,\widehat q)
+
r_{\mathrm{bc}}(\widehat u,\widehat q)
+
\delta
\right)^p .
\]
Since \(p\in(0,1]\) and all residual terms are nonnegative, substituting the upper bound
\[
r_{\mathrm{pde},\nu}(\widehat u,\widehat q)
\le
\widehat r_{\mathrm{pde},\zeta}(\widehat u,\widehat q)
\]
yields
\[
\|\widehat q-q^\dagger\|_{\mathcal Q}
\le
C_{\mathrm{stab}}
\left(
r_{\mathrm{data}}(\widehat u)
+
\widehat r_{\mathrm{pde},\zeta}(\widehat u,\widehat q)
+
r_{\mathrm{bc}}(\widehat u,\widehat q)
+
\delta
\right)^p .
\]
The right-hand side is \(R_{\delta,\zeta}^{\mathrm{stoch}}(\widehat u,\widehat q)\), proving \eqref{eq:stochastic-containment}.
\end{proof}

\subsection{High-probability no-harm selection}

The stochastic radius can be used directly as the certificate rule \(\mathfrak R\) in the no-harm selector. This gives a high-probability version of the no-harm guarantee.

\begin{corollary}[High-probability no-harm selection]
\label{cor:stochastic-no-harm-selection}
Let
\[
A_{\mathrm{base}}(y^\delta)
=
(u_{\mathrm{base}},q_{\mathrm{base}})
\]
and
\[
A_{\mathrm{learn}}(y^\delta)
=
(\widehat u,\widehat q)
\]
be fixed before drawing the independent validation collocation sets used to estimate their stochastic certificates. Suppose the assumptions of Corollary~\ref{cor:stochastic-containment} hold for both reconstructions, with confidence levels \(\zeta_{\mathrm{base}}\) and \(\zeta_{\mathrm{learn}}\). Use
\begin{equation}
\mathfrak R(A_{\mathrm{base}})
=
R_{\delta,\zeta_{\mathrm{base}}}^{\mathrm{stoch}}(u_{\mathrm{base}},q_{\mathrm{base}})
\label{eq:baseline-stoch-certificate}
\end{equation}
and
\begin{equation}
\mathfrak R(A_{\mathrm{learn}})
=
R_{\delta,\zeta_{\mathrm{learn}}}^{\mathrm{stoch}}(\widehat u,\widehat q)
\label{eq:learned-stoch-certificate}
\end{equation}
inside the no-harm selector. Then, with probability at least
\[
1-\zeta_{\mathrm{base}}-\zeta_{\mathrm{learn}},
\]
the selected output satisfies
\begin{equation}
\|q_{\mathrm{safe}}-q^\dagger\|_{\mathcal Q}
\le
R_{\delta,\zeta_{\mathrm{base}}}^{\mathrm{stoch}}(u_{\mathrm{base}},q_{\mathrm{base}})
+
\varepsilon_{\mathrm{safe}}.
\label{eq:stochastic-no-harm-selection}
\end{equation}
If the learned reconstruction is rejected, then with probability at least \(1-\zeta_{\mathrm{base}}\),
\begin{equation}
\|q_{\mathrm{safe}}-q^\dagger\|_{\mathcal Q}
\le
R_{\delta,\zeta_{\mathrm{base}}}^{\mathrm{stoch}}(u_{\mathrm{base}},q_{\mathrm{base}}).
\label{eq:stochastic-rejected-baseline}
\end{equation}
\end{corollary}

\begin{proof}
By Corollary~\ref{cor:stochastic-containment}, the baseline stochastic certificate is valid with probability at least \(1-\zeta_{\mathrm{base}}\), and the learned stochastic certificate is valid with probability at least \(1-\zeta_{\mathrm{learn}}\). By the union bound, both certificates are valid simultaneously with probability at least
\[
1-\zeta_{\mathrm{base}}-\zeta_{\mathrm{learn}}.
\]
On this simultaneous-validity event, Theorem~\ref{thm:no-harm-principle} applies with the certificate choices \eqref{eq:baseline-stoch-certificate} and \eqref{eq:learned-stoch-certificate}. Therefore,
\[
\|q_{\mathrm{safe}}-q^\dagger\|_{\mathcal Q}
\le
R_{\delta,\zeta_{\mathrm{base}}}^{\mathrm{stoch}}(u_{\mathrm{base}},q_{\mathrm{base}})
+
\varepsilon_{\mathrm{safe}},
\]
which proves \eqref{eq:stochastic-no-harm-selection}. If the learned reconstruction is rejected, then
\[
(u_{\mathrm{safe}},q_{\mathrm{safe}})
=
(u_{\mathrm{base}},q_{\mathrm{base}}).
\]
In that case, only the baseline stochastic certificate is needed, and Corollary~\ref{cor:stochastic-containment} applied to the baseline gives \eqref{eq:stochastic-rejected-baseline} with probability at least \(1-\zeta_{\mathrm{base}}\).
\end{proof}

The role of this result is practical. It shows how stochastic residual validation can be combined with the no-harm principle. A learned reconstruction is not selected because its training residual is small; it is selected only if its independently validated stochastic certificate is no worse than the corresponding baseline certificate, up to the prescribed tolerance.


\section{Optimization residual and approximate minimizers}
\label{sec:optimization}

A physics-informed reconstruction is normally the output of a finite, nonconvex training process. Data and physics residuals may be small because of the chosen sampling, scaling, or loss weights, while the optimizer is still far from stationarity. This is why the certificate includes an optimization residual. Stationarity residuals, projected-gradient norms, KKT residuals, and subdifferential distances are standard tools for assessing approximate solutions in numerical optimization \cite{NocedalWright2006,Bertsekas1999,RockafellarWets1998}.

Let \(\theta\) collect the trainable parameters of \(u_\theta\) and \(q_\theta\). For a differentiable training loss, define
\begin{equation}
r_{\mathrm{opt}}(\theta)
=
\|\nabla_\theta\mathcal L(\theta)\|.
\label{eq:opt-residual-repeat}
\end{equation}
For constrained formulations, \(r_{\mathrm{opt}}\) may instead be a projected-gradient norm. For nonsmooth formulations, it may be a subdifferential distance.

\begin{assumption}[Local optimization error bound]
\label{ass:local-error-bound}
There exist a neighbourhood \(\mathcal V\) of the computed parameter \(\widehat\theta\), a set \(\Theta^\star\) of stationary parameters satisfying the residual constraints, and a constant \(C_{\mathrm{opt}}>0\) such that
\begin{equation}
\dist(\theta,\Theta^\star)
\le
C_{\mathrm{opt}}\|\nabla_\theta\mathcal L(\theta)\|
\quad
\text{for all } \theta\in\mathcal V .
\label{eq:local-error-bound}
\end{equation}
\end{assumption}

Local error bounds of the form \eqref{eq:local-error-bound} are common in variational analysis and optimization. They hold for several structured problems under regularity conditions, and they are closely related to metric subregularity and Kurdyka--{\L}ojasiewicz-type inequalities \cite{LuoTseng1993,BolteSabachTeboulle2014,RockafellarWets1998}. Here the assumption is used only as a local diagnostic: if the optimization residual is large, the reconstruction should not be treated as fully certified even when the data and PDE residuals appear small.

\begin{proposition}[Optimization contribution to the certificate]
\label{prop:optimization-contribution}
Under Assumption~\ref{ass:local-error-bound}, the distance from the computed parameter \(\widehat\theta\) to the local stationary set satisfies
\begin{equation}
\dist(\widehat\theta,\Theta^\star)
\le
C_{\mathrm{opt}} r_{\mathrm{opt}}(\widehat\theta).
\label{eq:opt-distance-bound}
\end{equation}
Consequently, the term \(\alpha_{\mathrm{opt}}r_{\mathrm{opt}}(\widehat\theta)\) in \eqref{eq:total-certificate} records the part of the reliability certificate attributable to unfinished or inaccurate optimization.
\end{proposition}

\begin{proof}
The bound \eqref{eq:opt-distance-bound} is Assumption~\ref{ass:local-error-bound} evaluated at \(\theta=\widehat\theta\). The second statement follows from the definition of the total certificate \eqref{eq:total-certificate}.
\end{proof}

\subsection{Practical stopping rule}

The neural training process is treated as certifiable only when the total residual falls below a user-prescribed tolerance:
\begin{equation}
r_{\mathrm{data}}(\widehat u)
+
\alpha_{\mathrm{pde}} r_{\mathrm{pde}}(\widehat u,\widehat q)
+
\alpha_{\mathrm{bc}} r_{\mathrm{bc}}(\widehat u,\widehat q)
+
\alpha_{\mathrm{opt}} r_{\mathrm{opt}}(\widehat\theta)
\le
\tau_{\mathrm{cert}} .
\label{eq:stopping-rule}
\end{equation}
This rule prevents a reconstruction with small data or physics residuals but a large optimization residual from being reported as fully reliable. In implementation, the four residual components should also be reported separately, since a large certificate may be caused by data mismatch, physics violation, boundary error, or unfinished optimization.

The no-harm wrapper is summarized in Algorithm~\ref{alg:no-harm-piil}. The numerical validation below evaluates this same certification-and-selection rule in controlled inverse-problem settings.

\section{Numerical validation}
\label{sec:numerical-validation}

\subsection{Validation protocol}
\label{subsec:validation-protocol}

The numerical study tests whether the proposed certificate behaves as a reliability mechanism rather than as a post hoc accuracy score. The experiments address five questions. Does the residual-calibrated radius track reconstruction error in controlled inverse problems? Does the certified uncertainty set contain the true parameter when the ground truth is known? Does the no-harm rule accept learned reconstructions when their certified radius improves on the baseline? Does it reject learned reconstructions under noise shift, model mismatch, sparse data, hallucinated structure, or unfinished optimization? Does the same certification principle remain meaningful across PDE, tomography, and coefficient-identification examples?

All experiments are finite-dimensional and reproducible. They are designed to validate the certificate and the no-harm decision rule, not to claim that a neural architecture is necessary for the test. Each experiment compares a conservative baseline reconstruction with one or more learned or physics-informed candidate reconstructions. The baseline is used as the reference object in the no-harm rule. The learned candidates are not judged only by relative error or visual quality. They are judged by their data residual, physics residual, boundary or initial-condition residual, optimization residual, residual-calibrated radius, and no-harm decision.

Unless otherwise stated, the stability exponent is \(p=1\), the no-harm tolerance is \(\varepsilon_{\mathrm{safe}}=0\), and the boundary residual is zero because the discrete solvers impose the homogeneous Dirichlet boundary condition directly. All random perturbations and noise draws use a fixed random seed. The theoretical certified radius is defined in \eqref{eq:uncertainty-radius}. The numerical no-harm decisions use the operational radius \(R_{\delta}^{\mathrm{op}}\) in \eqref{eq:operational-radius}, with experiment-specific residual weights reported in Table~\ref{tab:validation-setup}. The conditional stability constant \(C_{\mathrm{stab}}\) is computed on the finite-dimensional admissible subspace as the reciprocal of the smallest singular value of the relevant observation-to-parameter map, or of the local finite-difference Jacobian in the nonlinear coefficient-identification example. These constants are finite-dimensional stability or conditioning constants for the stated discretized admissible spaces. They are used to make the no-harm decision transparent and reproducible, and should not be interpreted as sharp global stability constants for the corresponding infinite-dimensional inverse problems.

The reported metrics are the relative reconstruction error,
\[
\frac{\|\widehat q-q^\dagger\|}{\|q^\dagger\|},
\]
the coefficient or parameter error in the admissible coordinates, the data residual, the physics residual, the boundary or initial-condition residual, the optimization residual, the operational certified radius \(R_{\delta}^{\mathrm{op}}\), the radius sharpness ratio
\[
\frac{R_{\delta}^{\mathrm{op}}}{\|\widehat q-q^\dagger\|},
\]
the certificate coverage indicator, and the no-harm accept or reject decision. For imaging and tomography examples, visual reconstructions are also reported, but the no-harm rule itself is based on the certified radius, not on visual appearance. The certificate coverage indicator is a hindsight validation diagnostic used only when the exact \(q^\dagger\) is known. It equals one when the certified radius is at least as large as the observed reconstruction error, and zero otherwise. The radius sharpness ratio measures how conservative the certificate is: values closer to one indicate a tighter certificate, while large values indicate conservative but valid certification. Neither the coverage indicator nor the hindsight reconstruction error is used by the no-harm selector in deployment.

In the finite-dimensional experiments, data and parameter errors are computed using the discrete Euclidean norm on the corresponding observation, coefficient, grid, or residual vector. Relative errors are normalized by the Euclidean norm of the exact coefficient or field. For the Poisson source and elliptic coefficient experiments, the reported physics residual is root-mean-square (RMS)-scaled by the square root of the number of grid points, matching the implementation. For inverse heat and tomography, the physics residual is zero in the reported certificate because the tested candidates are evaluated through the discrete forward model or projection model used in the experiment. The same norm and scaling convention is used for the baseline and learned candidates in every no-harm comparison. This is important because the no-harm decision compares certified radii, so a different norm or scaling for the two candidates would invalidate the comparison.

\begin{table}[!ht]
\centering
\scriptsize
\caption{Numerical setup and parameter choices used in the validation experiments. Here \(n\) is the one-dimensional state or grid size when applicable, \(N\) is the tomography image width, \(k\) is the number of admissible basis functions, \(m\) is the number of observations, \(\lambda\) is the regularization parameter used in the baseline solve, and \(\varepsilon_{\mathrm{safe}}\) is the no-harm tolerance. A dash means that the quantity is not applicable to that experiment.}
\label{tab:validation-setup}
\resizebox{\textwidth}{!}{%
\begin{tabular}{lcccccccccc}
	\toprule
	\textbf{Experiment} & \(\boldsymbol{n}\) & \(\boldsymbol{N}\) & \(\boldsymbol{k}\) & \(\boldsymbol{m}\) & \textbf{Noise scale} & \(\boldsymbol{\lambda}\) & \(\boldsymbol{\alpha_{\mathrm{pde}}}\) & \(\boldsymbol{\alpha_{\mathrm{opt}}}\) & \(\boldsymbol{\varepsilon_{\mathrm{safe}}}\) & \textbf{Additional settings} \\
	\midrule
	Poisson source & 120 & -- & 10 & 35 & \(0.02\|\mathcal H u^\dagger\|/\sqrt m\) & \(10^{-5}\) & 0.05 & 0.01 & 0 & spread observations \\
	Inverse heat & 120 & -- & 8 & 45 & \(0.015\|\mathcal H u_T^\dagger\|/\sqrt m\) & \(10^{-4}\) & 0 & 0.005 & 0 & \(T\in\{0.02,0.08,0.16\}\), \(\kappa=0.004\) \\
	Limited-angle tomography & -- & 28 & 36 & 420 & \(0.01\|\mathcal R_\Omega q^\dagger\|/\sqrt m\) & \(10^{-3}\) & 0 & 0.001 & 0 & 15 angles in \([-50^\circ,50^\circ]\) \\
	Elliptic coefficient & 90 & -- & 6 & 30 & \(0.01\|\mathcal H u^\dagger\|/\sqrt m\) & \(5\times10^{-3}\) & 0.01 & 0.001 & 0 & forcing \(f(x)=1+0.5\sin(2\pi x)\) \\
	Stochastic residual sweep & 2000 & -- & -- & -- & -- & -- & -- & -- & -- & \(M\in\{20,50,100,200,500,1000,2000\}\), 250 repetitions, \(\zeta=0.05\) \\
	\bottomrule
\end{tabular}%
}
\end{table}

The learned candidates are controlled physics-informed candidate reconstructions used to stress-test the certificate. Some are constructed to be close to the truth, while others deliberately include shifted identifiable modes, hallucinated high-frequency structure, or state--source inconsistency. This design tests whether the no-harm selector accepts candidates supported by the certificate and rejects candidates whose residual-calibrated evidence is weaker than the baseline. The validation therefore evaluates the certification-and-selection mechanism, not the superiority of a particular neural-network architecture.

The first validation setting is inverse source recovery for the Poisson equation,
\begin{equation}
-\Delta u=q
\quad \text{in } \Omega,
\qquad
u=0
\quad \text{on } \partial\Omega .
\label{eq:poisson-source}
\end{equation}
The one-dimensional domain is discretized with \(n=120\) interior points. The unknown source is represented as \(q=Bc\), where \(B\) is a normalized sine basis with \(k=10\) columns and \(c\in\mathbb R^{10}\) is the admissible coefficient vector. We use \(m=35\) spread observations of the state. The finite-difference Poisson matrix maps the source to the state, and the admissible observation map is \(F=\mathcal H K^{-1}B\), where \(B\) is the source basis and \(K\) is the discrete Poisson operator. The baseline is the ridge solution with \(\lambda=10^{-5}\). The good learned candidate is generated by perturbing the true coefficient vector by \(0.035\) times a standard normal vector. The shifted learned candidate adds the structured perturbation \((0.65,-0.45,0.35)\) to three identifiable coefficient modes. The unfinished PINN-like candidate perturbs the coefficients by \(0.06\) times a standard normal vector and adds an inconsistent state perturbation \(0.06\sin(15\pi x)\). The radius uses \(\alpha_{\mathrm{pde}}=0.05\) and \(\alpha_{\mathrm{opt}}=0.01\).

The second validation setting is inverse heat reconstruction,
\begin{equation}
\partial_t u-\kappa\Delta u=0
\quad \text{in } \Omega\times(0,T],
\qquad
u(\cdot,0)=q .
\label{eq:inverse-heat}
\end{equation}
The unknown initial condition is represented as \(q=Bc\), where \(B\) is a normalized sine basis with \(k=8\) columns and \(c\in\mathbb R^8\) is the admissible coefficient vector on a grid with \(n=120\) interior points. The diffusivity is \(\kappa=0.004\), and the final observation times are \(T=0.02\), \(T=0.08\), and \(T=0.16\). For each \(T\), \(m=45\) spread observations are taken from the propagated state \(u_T\). The heat propagator is computed from the eigendecomposition of the finite-difference Laplacian. The baseline is a ridge reconstruction with \(\lambda=10^{-4}\). The good learned candidate perturbs the true coefficient vector by \(0.04\) times a standard normal vector. The shifted learned candidate adds \((0.65,-0.50,0.35,-0.25)\) to four middle coefficient modes. The hallucinated high-frequency candidate adds \(0.15\sin(30\pi x)\) to the good learned reconstruction. The radius uses \(\alpha_{\mathrm{opt}}=0.005\), with zero physics residual because the candidate states are generated through the discrete heat propagator.

The third validation setting is limited-angle tomography,
\begin{equation}
y^\delta=\mathcal R_\Omega q^\dagger+e ,
\label{eq:limited-angle-tomography}
\end{equation}
where \(\mathcal R_\Omega\) denotes the finite-dimensional restricted Radon-like projection operator used in the experiment, with projections taken only over the available limited-angle set. The image has size \(28\times28\), and the admissible representation uses a two-dimensional cosine basis with \(6\times6=36\) basis functions. The projection data are generated from 15 angles uniformly spaced in \([-50^\circ,50^\circ]\), giving \(420\) measurements. The baseline is a ridge reconstruction with \(\lambda=10^{-3}\). The good learned candidate perturbs the admissible coefficient vector by \(0.03\) times a standard normal vector. The hallucinated learned candidate adds the structured perturbation \((2.0,-1.8,1.4,-1.1,0.8)\) to selected coefficient modes. The radius uses \(\alpha_{\mathrm{opt}}=0.001\). This experiment tests whether the no-harm rule can reject visually plausible learned structure that is weakly supported by limited-angle projection data.

The fourth validation setting is geophysical-style elliptic coefficient identification,
\begin{equation}
-\nabla\cdot(a(x)\nabla u)=f
\quad \text{in } \Omega,
\qquad
u=0
\quad \text{on } \partial\Omega .
\label{eq:elliptic-coefficient}
\end{equation}
The coefficient is parameterized as \(a(x)=\exp(Bc)(x)\), where \(B\) is a normalized sine basis with \(k=6\) columns. The grid has \(n=90\) interior points, and \(m=30\) spread observations of the state are used. The forcing is
\[
f(x)=1+0.5\sin(2\pi x).
\]
The baseline is obtained by nonlinear least squares from the zero coefficient vector, using ridge weight \(\lambda=5\times10^{-3}\), maximum function evaluations \(300\), and solver stopping tolerances \(10^{-10}\) for the step-size tolerance \texttt{xtol}, cost-change tolerance \texttt{ftol}, and gradient tolerance \texttt{gtol}. The local stability constant is computed from a finite-difference Jacobian at the true coefficient vector using step size \(10^{-5}\). The good learned candidate perturbs the true coefficient vector by \(0.025\) times a standard normal vector. The shifted learned candidate adds \((0.40,-0.25,0.20,-0.18,0.10,-0.08)\) to the coefficient vector. The radius uses \(\alpha_{\mathrm{pde}}=0.01\) and \(\alpha_{\mathrm{opt}}=0.001\).

The stochastic residual experiment evaluates the effect of estimating the physics residual from independent validation collocation points. The residual field is defined on \(n=2000\) grid points by
\[
r(x)
=
0.3\sin(6\pi x)
+
1.5\exp\!\left(-\frac{(x-0.72)^2}{0.002}\right)
+
0.8\exp\!\left(-\frac{(x-0.22)^2}{0.0008}\right).
\]
The true mean squared residual (MSR) is \(0.225193\), where MSR denotes the squared physics residual averaged over the validation distribution. The residual-envelope constant in Theorem~\ref{thm:high-prob-residual} is
\[
B_{\mathrm{pde}}=3.084186 .
\]
For each sample size
\[
M\in\{20,50,100,200,500,1000,2000\},
\]
the experiment draws \(250\) independent validation samples with replacement, computes the empirical squared residual, and evaluates the high-probability upper bound at confidence parameter \(\zeta=0.05\). This test checks whether the stochastic certificate is conservative and whether the upper bound tightens as the number of independent validation collocation points increases.

\subsection{Sufficiency map for no-harm selection}
\label{subsec:sufficiency-map}

The preceding theory gives a decision rule for using a learned physics-informed inverse reconstruction. The remaining question is operational: under what conditions is the learned reconstruction certified enough to replace the baseline? To answer this, we run a controlled sufficiency sweep for the Poisson inverse source problem. The sweep is not intended to introduce a new benchmark. Its purpose is to map the no-harm decision boundary across observation density, noise level, learned-reconstruction perturbation, and physics inconsistency.

The central quantity is the certificate ratio
\begin{equation}
\Gamma
=
\frac{R_{\mathrm{learn}}}{R_{\mathrm{base}}}.
\label{eq:certificate-ratio}
\end{equation}
With the strict no-harm choice \(\varepsilon_{\mathrm{safe}}=0\), the learned reconstruction is selected exactly when
\begin{equation}
\Gamma \le 1.
\label{eq:sufficiency-threshold}
\end{equation}
Thus the practical sufficiency condition is
\begin{equation}
\boxed{
\text{physics-informed inverse learning is sufficient to replace the baseline only when }
R_{\mathrm{learn}}\le R_{\mathrm{base}}+\varepsilon_{\mathrm{safe}}.
}
\label{eq:boxed-sufficiency-condition}
\end{equation}
This is the threshold tested in this subsection. The hindsight error ratio
\begin{equation}
E_{\mathrm{ratio}}
=
\frac{\|q_{\mathrm{learn}}-q^\dagger\|_{\mathcal Q}}
{\|q_{\mathrm{base}}-q^\dagger\|_{\mathcal Q}}
\label{eq:error-ratio}
\end{equation}
is reported only for validation, because \(q^\dagger\) is unavailable in real inverse problems. The no-harm selector does not use \eqref{eq:error-ratio}. It uses the certificate ratio \eqref{eq:certificate-ratio}.

Table~\ref{tab:sufficiency-sweep-design} reports the sweep design. The experiment uses a one-dimensional Poisson inverse source problem with \(n=120\) state grid points and \(k=10\) admissible source-basis coefficients. The number of observations, noise level, learned coefficient perturbation, and state-source mismatch are varied systematically. Each regime is repeated \(25\) times, giving \(840\) regimes and \(21{,}000\) trials.

Unlike the representative Poisson experiment in Table~\ref{tab:validation-setup}, the sufficiency sweep uses randomly selected observation locations for each trial. For a noise fraction \(\eta\), the noise standard deviation is scaled as \(\eta\|\mathcal H u^\dagger\|/\sqrt m\), matching the validation code. This makes the sweep a stress test over observation density, random sensing geometry, noise level, learned coefficient perturbation, and state--source inconsistency.

\begin{table}[!ht]
\centering
\caption{Design of the no-harm sufficiency sweep. The sweep varies observation density, noise, learned coefficient perturbation, and physics inconsistency in a controlled Poisson inverse source problem.}
\label{tab:sufficiency-sweep-design}
\begin{tabular}{ll}
	\toprule
	\textbf{Quantity} & \textbf{Value} \\
	\midrule
	State grid size & \(n=120\) \\
	Admissible source basis size & \(k=10\) \\
	Observation counts & \(m\in\{10,15,20,25,35,50,70\}\) \\
	Noise fractions & \(\eta\in\{0,0.01,0.02,0.05,0.10\}\) \\
	Learned coefficient perturbation scales & \(\sigma_{\mathrm{learn}}\in\{0,0.02,0.05,0.10,0.20,0.40\}\) \\
	State-source mismatch amplitudes & \(\mu\in\{0,0.02,0.05,0.10\}\) \\
	Replications per regime & \(25\) \\
	Total regimes & \(840\) \\
	Total trials & \(21{,}000\) \\
	Baseline regularization & \(\lambda=10^{-5}\) \\
	No-harm tolerance & \(\varepsilon_{\mathrm{safe}}=0\) \\
	Operational residual weights & \(\alpha_{\mathrm{pde}}=0.05\), \(\alpha_{\mathrm{opt}}=0.01\) \\
	\bottomrule
\end{tabular}
\end{table}

Figure~\ref{fig:suff-cert-boundary} shows the main decision boundary. The vertical line \(\Gamma=1\) separates the certificate-sufficient region from the fallback region. Points to the left of this line are learned reconstructions whose certified radius is no worse than the baseline radius. Points to the right are learned reconstructions whose certificate is weaker than the baseline certificate. The horizontal line \(E_{\mathrm{ratio}}=1\) is shown only for validation, since true error is unavailable in deployment.

\begin{figure}[!ht]
\centering
\includegraphics[width=0.6\textwidth]{certificate_boundary_scatter.pdf}
\caption{No-harm sufficiency boundary. The learned reconstruction is selected only when \(R_{\mathrm{learn}}/R_{\mathrm{base}}\le1\). The horizontal error-ratio line is a hindsight diagnostic and is not used by the selector. The figure shows the central distinction: certificate dominance, not visual quality or hindsight error, is the operational threshold for replacing the baseline.}
\label{fig:suff-cert-boundary}
\end{figure}

Table~\ref{tab:sufficiency-aggregate-outcomes} summarizes the full sweep. Out of \(21{,}000\) trials, the no-harm rule selected the learned reconstruction in \(1{,}235\) cases, or \(5.88\%\) of all trials. Among selected learned reconstructions, \(1{,}221\) were true improvements in hindsight, corresponding to \(98.87\%\) of selections. Unsafe selections were rare: \(14\) cases, or \(0.07\%\) of all trials and \(1.13\%\) of selected cases. The method also produced many hindsight false rejections, \(11{,}499\) cases, or \(54.76\%\) of all trials. This is expected for a conservative no-harm rule: the selector is not an oracle for the unknown error, and it rejects learned reconstructions when their available certificate does not justify replacing the baseline.

\begin{table}[!ht]
\centering
\caption{Aggregate outcomes of the no-harm sufficiency sweep. A safe improvement means that the learned reconstruction was selected and had lower hindsight error than the baseline. An unsafe selection means that the learned reconstruction was selected but had higher hindsight error. A false rejection means that the learned reconstruction had lower hindsight error but was rejected because its certificate was weaker than the baseline certificate.}
\label{tab:sufficiency-aggregate-outcomes}
\begin{tabular}{lrr}
	\toprule
	\textbf{Quantity} & \textbf{Count} & \textbf{Rate} \\
	\midrule
	Total trials & \(21{,}000\) & \(100.00\%\) \\
	Selected by no-harm rule & \(1{,}235\) & \(5.88\%\) of all trials \\
	Safe improvements & \(1{,}221\) & \(5.81\%\) of all trials, \(98.87\%\) of selections \\
	Unsafe selections & \(14\) & \(0.07\%\) of all trials, \(1.13\%\) of selections \\
	Hindsight false rejections & \(11{,}499\) & \(54.76\%\) of all trials \\
	Certificate-sufficient regimes & \(41\) & \(4.88\%\) of \(840\) regimes \\
	Fallback-required regimes & \(799\) & \(95.12\%\) of \(840\) regimes \\
	\bottomrule
\end{tabular}
\end{table}

Table~\ref{tab:sufficiency-regime-summary} separates regimes by the median certificate ratio. In certificate-sufficient regimes, the median certificate ratio is \(0.6425\), the median hindsight error ratio is \(0.0774\), and the learned reconstruction is selected in \(88.68\%\) of trials. In fallback-required regimes, the median certificate ratio is \(1638.9831\), the median hindsight error ratio is \(1.2840\), and the learned reconstruction is selected in only \(1.63\%\) of trials. This confirms that the certificate ratio is acting as the intended operational boundary.

\begin{table}[!ht]
\centering
\caption{Regime-level summary of the sufficiency sweep. Certificate-sufficient regimes satisfy median \(R_{\mathrm{learn}}/R_{\mathrm{base}}\le1\). Fallback-required regimes have median \(R_{\mathrm{learn}}/R_{\mathrm{base}}>1\).}
\label{tab:sufficiency-regime-summary}
\resizebox{\textwidth}{!}{%
\begin{tabular}{lrrrrrrr}
	\toprule
	\textbf{Regime class} & \textbf{Regimes} & \(\boldsymbol{\mathrm{median}\ \Gamma}\) & \(\boldsymbol{\mathrm{median}\ E_{\mathrm{ratio}}}\) & \textbf{Accept rate} & \textbf{Learned better rate} & \textbf{Unsafe selection rate} & \textbf{False rejection rate} \\
	\midrule
	Certificate-sufficient & \(41\) & \(0.6425\) & \(0.0774\) & \(88.68\%\) & \(100.00\%\) & \(0.00\%\) & \(11.32\%\) \\
	Fallback-required & \(799\) & \(1638.9831\) & \(1.2840\) & \(1.63\%\) & \(58.55\%\) & \(0.07\%\) & \(56.99\%\) \\
	\bottomrule
\end{tabular}}
\end{table}

Table~\ref{tab:mismatch-regime-summary} shows why physics consistency matters. When the state-source mismatch amplitude is zero, the no-harm selector accepts learned reconstructions in \(23.52\%\) of trials, with median certificate ratio \(1.2464\). Once mismatch is introduced, the selector becomes sharply conservative: the acceptance rate is \(0.00\%\) for \(\mu=0.02\), \(\mu=0.05\), and \(\mu=0.10\). The median certificate ratio rises from \(633.6613\) to \(3121.7369\) as the mismatch amplitude increases. This is the behaviour expected from a physics-informed certificate: a learned reconstruction is not enough merely because its coefficient error may be low; it must also be consistent with the governing equation and the observation evidence.

\begin{table}[!ht]
\centering
\caption{Effect of state-source mismatch on no-harm selection. The mismatch amplitude \(\mu\) measures physics inconsistency in the learned candidate. Even when the learned reconstruction is better in hindsight, mismatch can make its certificate too weak to replace the baseline.}
\label{tab:mismatch-regime-summary}
\begin{tabular}{rrrrrr}
	\toprule
	\(\boldsymbol{\mu}\) & \textbf{Accept rate} & \(\boldsymbol{\mathrm{median}\ \Gamma}\) & \(\boldsymbol{\mathrm{median}\ E_{\mathrm{ratio}}}\) & \textbf{Unsafe selection rate} & \textbf{False rejection rate} \\
	\midrule
	\(0.00\) & \(23.52\%\) & \(1.2464\) & \(0.6512\) & \(0.27\%\) & \(36.90\%\) \\
	\(0.02\) & \(0.00\%\) & \(633.6613\) & \(0.6387\) & \(0.00\%\) & \(61.12\%\) \\
	\(0.05\) & \(0.00\%\) & \(1569.9048\) & \(0.6418\) & \(0.00\%\) & \(60.13\%\) \\
	\(0.10\) & \(0.00\%\) & \(3121.7369\) & \(0.6337\) & \(0.00\%\) & \(60.88\%\) \\
	\bottomrule
\end{tabular}
\end{table}

Figures~\ref{fig:suff-acceptance-heatmap}--\ref{fig:suff-false-rejection-heatmap} give the corresponding heatmap view for the representative setting \(\eta=0.02\) and \(\mu=0.05\). Figure~\ref{fig:suff-acceptance-heatmap} shows that no learned candidate is accepted in this mismatched setting. Figure~\ref{fig:suff-unsafe-heatmap} shows zero unsafe acceptance over the same grid. Figure~\ref{fig:suff-false-rejection-heatmap} shows the cost of this conservativeness: accurate learned candidates may still be rejected when their certificate is weakened by physics inconsistency. This is not a defect in the rule. It is the intended no-harm behaviour under uncertainty.

\begin{figure}[!ht]
\begin{minipage}{0.48\textwidth}
\centering
\includegraphics[width=\textwidth]{acceptance_rate_heatmap.pdf}
\caption{Acceptance-rate heatmap for the representative mismatched setting \(\eta=0.02\) and \(\mu=0.05\). Across the displayed observation counts and learned perturbation scales, the no-harm selector does not accept the learned reconstruction because physics inconsistency makes the learned certificate weaker than the baseline certificate.}
\label{fig:suff-acceptance-heatmap}
\end{minipage}
\hfill
\begin{minipage}{0.48\textwidth}
\centering
\includegraphics[width=\textwidth]{unsafe_acceptance_heatmap.pdf}
\caption{Unsafe-acceptance heatmap for the representative mismatched setting \(\eta=0.02\) and \(\mu=0.05\). The unsafe acceptance rate is zero over the displayed grid, showing that the no-harm selector blocks learned candidates whose certificates do not justify replacing the baseline.}
\label{fig:suff-unsafe-heatmap}
\end{minipage}
\end{figure}

\begin{figure}[!ht]
\begin{minipage}{0.48\textwidth}
\centering
\includegraphics[width=\textwidth]{false_rejection_heatmap.pdf}
\caption{False-rejection heatmap for the representative mismatched setting \(\eta=0.02\) and \(\mu=0.05\). False rejection is a hindsight diagnostic: it occurs when a learned reconstruction is more accurate than the baseline but is rejected because the certificate is weaker. This illustrates the conservative nature of no-harm selection.}
\label{fig:suff-false-rejection-heatmap}
\end{minipage}
\hfill
\begin{minipage}{0.48\textwidth}
\centering
\includegraphics[width=\textwidth]{certificate_ratio_by_data_noise.pdf}
\caption{Median certificate ratio \(R_{\mathrm{learn}}/R_{\mathrm{base}}\) over observation count and noise level for \(\sigma_{\mathrm{learn}}=0.05\) and \(\mu=0.02\). Values above one lie outside the certificate-sufficient region. The figure shows that increased observation density can reduce the ratio, but physics inconsistency can still keep the learned reconstruction outside the no-harm acceptance region.}
\label{fig:suff-cert-ratio-data-noise}
\end{minipage}
\end{figure}

Figure~\ref{fig:suff-cert-ratio-data-noise} shows the median certificate ratio as a function of observation count and noise level for a representative learned perturbation scale \(\sigma_{\mathrm{learn}}=0.05\) and mismatch amplitude \(\mu=0.02\). The median certificate ratio ranges from \(89.48\) to \(2065.60\), so all displayed regimes are outside the certificate-dominance region. Increasing the number of observations reduces the ratio in many noisy regimes, but it does not by itself overcome physics inconsistency. This supports the central message: more data can help the certificate, but a learned inverse output still needs residual-calibrated certificate dominance before it is allowed to replace the baseline.

The sufficiency sweep answers the threshold question directly. A physics-informed learned reconstruction is not sufficient because it is physics-informed, neural, visually plausible, or accurate in hindsight. It is sufficient to replace the baseline only when its residual-calibrated certificate dominates the baseline certificate. In the strict case used here, the operational threshold is \(R_{\mathrm{learn}}/R_{\mathrm{base}}\le1\). Below this threshold, the learned reconstruction is selected. Above it, the method falls back to the baseline. The subsequent validation examples test this same rule in Poisson source recovery, inverse heat reconstruction, limited-angle tomography, geophysical coefficient identification, and stochastic residual certification.

\subsection{Results and no-harm decisions}
\label{subsec:results-noharm-decisions}

The numerical results now test the same no-harm threshold in representative inverse-problem regimes after the sufficiency map in Section~\ref{subsec:sufficiency-map}. The experiments are not designed to show that a learned reconstruction is always superior. They test a narrower claim: a physics-informed learned reconstruction should be accepted when its residual-calibrated radius improves on the baseline, and rejected when its apparent accuracy is not supported by the data, physics, boundary, or optimization certificate.

In each experiment, we compute a baseline reconstruction, one or more learned reconstruction candidates, the residual components, the certified radius, the coverage indicator when the ground truth is known, and the no-harm decision. Full candidate-level results are reported in Appendix~\ref{app:full-validation-results}; the main text reports the key decision outcomes.

Figures~\ref{fig:reconstructions-poisson-heat}--\ref{fig:tomography-3d} report the visual reconstruction and diagnostic evidence used in the validation. Figure~\ref{fig:reconstructions-poisson-heat} shows the Poisson source and inverse heat examples. Figure~\ref{fig:reconstructions-geophys-residual} shows the geophysical coefficient example and the residual field used for stochastic residual certification. Figures~\ref{fig:tomography-reconstruction} and \ref{fig:tomography-3d} report the limited-angle tomography reconstruction comparison and the corresponding three-dimensional surface view.

\begin{figure}[!t]
\centering
\begin{minipage}{0.48\textwidth}
\centering
\includegraphics[width=\linewidth]{poisson_source_reconstructions.pdf}

\smallskip
\textbf{Left panel.} Poisson source recovery.
\end{minipage}
\hfill
\begin{minipage}{0.48\textwidth}
\centering
\includegraphics[width=\linewidth]{heat_initial_condition_panel.pdf}

\smallskip
\textbf{Right panel.} Inverse heat initial-condition reconstruction.
\end{minipage}
\caption{Representative PDE-based reconstructions for the Poisson and inverse heat experiments. The left panel shows source recovery in a controlled elliptic inverse problem. The right panel shows initial-condition recovery in a parabolic inverse problem, where smoothing makes the reconstruction less stable.}
\label{fig:reconstructions-poisson-heat}
\end{figure}


\begin{figure}[!t]
\centering
\begin{minipage}{0.48\textwidth}
\centering
\includegraphics[width=\linewidth]{geophysical_coefficient_reconstructions.pdf}

\smallskip
\textbf{Left panel.} Elliptic coefficient recovery.
\end{minipage}
\hfill
\begin{minipage}{0.48\textwidth}
\centering
\includegraphics[width=\linewidth]{stochastic_residual_field.pdf}

\smallskip
\textbf{Right panel.} Physics-residual field.
\end{minipage}
\caption{Coefficient-recovery and residual-field diagnostics. The left panel shows the geophysical-style elliptic coefficient-identification example, where the unknown enters the differential operator. The right panel shows the physics-residual field used to evaluate the stochastic residual certificate on independent validation collocation points.}
\label{fig:reconstructions-geophys-residual}
\end{figure}


\begin{figure}[!t]
\centering
\includegraphics[width=\linewidth]{tomography_reconstruction_panel.pdf}
\caption{Limited-angle tomography reconstruction comparison. The figure compares the baseline, learned, and no-harm safe outputs for the tomography example.}
\label{fig:tomography-reconstruction}
\end{figure}


\begin{figure}[!t]
\centering
\includegraphics[width=\linewidth]{tomography_3d_surface_truth_vs_safe.pdf}
\caption{Three-dimensional tomography surface comparison. The figure compares the true reconstruction target with the no-harm safe output in surface form.}
\label{fig:tomography-3d}
\end{figure}

Table~\ref{tab:main-noharm-summary} reports the main no-harm decisions. The Poisson source experiment accepts the good learned reconstruction because its certified radius is smaller than the baseline radius. The shifted learned reconstruction and unfinished PINN are rejected. In the inverse heat experiment, all representative learned candidates are rejected at \(T=0.08\), because their certified radii exceed the baseline radius. In limited-angle tomography, the learned-good candidate has a smaller relative error than the baseline, but it is still rejected because its certified radius is larger. This is an important stress test: the no-harm rule is not an error oracle; it is a reliability rule based on computable certificate quantities. In the geophysical coefficient experiment, the good learned coefficient is accepted, while the shifted learned coefficient is rejected.

\begin{table}[!ht]
\centering
\caption{Main no-harm decision outcomes. \(R_{\mathrm{base}}\) and \(R_{\mathrm{learn}}\) denote the operational certified radii used by the numerical no-harm rule. The learned candidate is selected only when \(R_{\mathrm{learn}}\le R_{\mathrm{base}}+\varepsilon_{\mathrm{safe}}\); otherwise the selected output falls back to the baseline.}
\label{tab:main-noharm-summary}
\resizebox{\linewidth}{!}{
\begin{tabular}{lclcccl}
	\toprule
	\textbf{Problem} & \(\boldsymbol{R_{\mathrm{base}}}\) & \textbf{Candidate} & \(\boldsymbol{R_{\mathrm{learn}}}\) & \textbf{Rel. error} & \textbf{Decision} & \textbf{Safe output} \\
	\midrule
	Poisson source & 5.375 & learned good & 4.930 & 0.0719 & accept & learned \\
	Poisson source & 5.375 & shifted learned & 19.28 & 0.5739 & reject & baseline \\
	Poisson source & 5.375 & unfinished PINN & 9045 & 0.1020 & reject & baseline \\
	Inverse heat \((T=0.08)\) & 0.0403 & learned good & 0.1181 & 0.0748 & reject & baseline \\
	Inverse heat \((T=0.08)\) & 0.0403 & hallucinated high freq. & 0.1612 & 0.9604 & reject & baseline \\
	Inverse heat \((T=0.08)\) & 0.0403 & shifted learned & 1.116 & 0.7599 & reject & baseline \\
	Tomography & 6.520 & learned good & 7.922 & 0.0190 & reject & baseline \\
	Tomography & 6.520 & hallucinated learned & 76.93 & 0.3204 & reject & baseline \\
	Geophysical coeff. & 6.048 & learned good & 4.524 & 0.0050 & accept & learned \\
	Geophysical coeff. & 6.048 & shifted learned & 11.01 & 0.0635 & reject & baseline \\
	\bottomrule
\end{tabular}}
\end{table}

Figures~\ref{fig:certificate-poisson-heat}--\ref{fig:certificate-coverage-decisions} report the certificate-level diagnostics. Figure~\ref{fig:certificate-poisson-heat} compares true error with the certified radius in the Poisson source experiment and shows how the inverse heat stability constant changes with the final observation time. Figure~\ref{fig:certificate-stochastic-aggregate} reports the stochastic residual sweep and the aggregate error-radius behaviour across experiments. Figure~\ref{fig:certificate-coverage-decisions} summarizes empirical coverage and no-harm decisions across the validation experiments.

\begin{figure}[!t]
\centering
\begin{minipage}{0.48\textwidth}
\centering
\includegraphics[width=\linewidth]{poisson_error_vs_radius.pdf}

\smallskip
\textbf{Left panel.} Poisson error versus certified radius.
\end{minipage}
\hfill
\begin{minipage}{0.48\textwidth}
\centering
\includegraphics[width=\linewidth]{heat_stability_constant_vs_time.pdf}

\smallskip
\textbf{Right panel.} Heat inverse stability constant.
\end{minipage}
\caption{Problem-specific certificate diagnostics. The left panel compares the true reconstruction error with the certified radius in the Poisson source experiment. The right panel shows how the stability constant used in the inverse heat certificate changes as the final observation time increases.}
\label{fig:certificate-poisson-heat}
\end{figure}

\begin{figure}[!t]
\centering
\begin{minipage}{0.48\textwidth}
\centering
\includegraphics[width=\linewidth]{stochastic_residual_sweep.pdf}

\smallskip
\textbf{Left panel.} Stochastic residual sweep.
\end{minipage}
\hfill
\begin{minipage}{0.48\textwidth}
\centering
\includegraphics[width=\linewidth]{aggregate_error_vs_radius.pdf}

\smallskip
\textbf{Right panel.} Aggregate error-radius behaviour.
\end{minipage}
\caption{Residual-sampling and aggregate calibration diagnostics. The left panel shows the stochastic residual sweep used to assess the high-probability residual certificate. The right panel summarizes the relationship between true reconstruction error and certified radius across the validation experiments.}
\label{fig:certificate-stochastic-aggregate}
\end{figure}

\begin{figure}[!t]
\centering
\begin{minipage}{0.48\textwidth}
\centering
\includegraphics[width=\linewidth]{aggregate_coverage_by_experiment.pdf}

\smallskip
\textbf{Left panel.} Coverage by experiment.
\end{minipage}
\hfill
\begin{minipage}{0.48\textwidth}
\centering
\includegraphics[width=\linewidth]{aggregate_noharm_decisions.pdf}

\smallskip
\textbf{Right panel.} No-harm decisions.
\end{minipage}
\caption{Aggregate decision diagnostics. The left panel reports empirical coverage of the certified uncertainty sets by experiment. The right panel summarizes the no-harm accept and reject decisions across the validation cases. Together, these diagnostics show that the residual-calibrated certificate is conservative in these tests and that the no-harm rule accepts learned reconstructions only when their certified radius improves on the baseline.}
\label{fig:certificate-coverage-decisions}
\end{figure}

Table~\ref{tab:stochastic-residual-main} reports the stochastic residual validation. The empirical squared residual fluctuates around the true mean squared residual, while the high-probability upper bound decreases as the number of independent validation collocation points increases. The observed coverage rate is \(1.000\) across the tested sample sizes. This supports the use of an independent residual-validation set in the stochastic certificate.


\begin{table}[!ht]
\centering
\caption{Stochastic residual validation. The true mean squared residual is \(0.2252\). The high-probability upper bound decreases with the number of independent validation collocation points \(M\).}
\label{tab:stochastic-residual-main}
\begin{tabular}{rrrrrrr}
	\toprule
	\(\boldsymbol{M}\) & \textbf{Reps.} & \textbf{True MSR} & \textbf{Empirical MSR} & \textbf{HP upper bound} & \textbf{Coverage} & \(\boldsymbol{\zeta}\) \\
	\midrule
	20 & 250 & 0.2252 & 0.2313 & 1.075 & 1.000 & 0.05 \\
	50 & 250 & 0.2252 & 0.2279 & 0.7618 & 1.000 & 0.05 \\
	100 & 250 & 0.2252 & 0.2209 & 0.5984 & 1.000 & 0.05 \\
	200 & 250 & 0.2252 & 0.2245 & 0.4914 & 1.000 & 0.05 \\
	500 & 250 & 0.2252 & 0.2238 & 0.3926 & 1.000 & 0.05 \\
	1000 & 250 & 0.2252 & 0.2243 & 0.3437 & 1.000 & 0.05 \\
	2000 & 250 & 0.2252 & 0.2237 & 0.3081 & 1.000 & 0.05 \\
	\bottomrule
\end{tabular}
\end{table}

The full numerical outputs are placed in Appendix~\ref{app:full-validation-results}. Appendix Table~\ref{tab:appendix-noharm-full} reports all no-harm decisions, including the inverse heat experiments at \(T=0.02\), \(T=0.08\), and \(T=0.16\). Appendix Table~\ref{tab:appendix-stability-constants} reports the stability constants used in the radius calculations. Appendix Table~\ref{tab:appendix-stochastic-sweep} repeats the stochastic residual sweep in full for reproducibility.

\section{Discussion}
\label{sec:discussion}

The numerical validation supports the intended role of the proposed framework. No-harm physics-informed inverse learning is not designed to make every learned reconstruction acceptable. It is a certification-and-selection principle: a learned reconstruction is allowed to replace a safer baseline only when its residual-calibrated certificate supports that replacement. The decisive condition is certificate dominance,
\[
R_{\mathrm{learn}}
\le
R_{\mathrm{base}}+\varepsilon_{\mathrm{safe}}.
\]
When this condition holds, the learned reconstruction is selected. When it fails, the method falls back to the baseline. This gives a clear operational threshold for using physics-informed learning in inverse problems: the learned reconstruction must be better supported by the available residual-calibrated evidence, not merely visually plausible, neural, or accurate in hindsight.

The sufficiency sweep makes this threshold explicit. With the strict no-harm choice \(\varepsilon_{\mathrm{safe}}=0\), the learned reconstruction is selected exactly in the certificate-dominance region \(R_{\mathrm{learn}}/R_{\mathrm{base}}\le1\). Across \(21{,}000\) trials, the no-harm selector accepted the learned reconstruction in \(1{,}235\) cases. Of these accepted cases, \(1{,}221\) were true improvements in hindsight. Unsafe selections were rare, with \(14\) cases. The sweep also produced many hindsight false rejections, where a learned reconstruction was more accurate than the baseline but was rejected because its certificate was weaker. This is not a defect of the rule. It is the conservative behaviour expected from a reliability principle. In real inverse problems, the true error is unavailable, so the method must act on the certificate available at reconstruction time.

The representative inverse-problem experiments show both sides of the rule. In the Poisson source and geophysical coefficient experiments, the good learned candidates are accepted because their certified radii are smaller than the baseline radii. Shifted, hallucinated, and unfinished learned candidates are rejected because their certified radii are larger. These outcomes show that the method is not biased against learned solvers. It accepts learned reconstructions when the certificate improves, and it blocks them when the residual-calibrated evidence is weaker than the baseline.

The limited-angle tomography result is especially important. The learned-good candidate has a smaller relative error than the baseline, but it is rejected because its certified radius is larger. This does not contradict the method. The no-harm rule is not an oracle for the unknown ground-truth error. It is a deployment-time reliability rule. In real tomography and imaging problems, the true object is not available, and visual plausibility is not a certificate. The relevant question is whether the reconstruction is supported by data consistency, model consistency, residual control, and the assumed stability estimate. The tomography result therefore shows the central point of the paper: a learned inverse reconstruction may be accurate in hindsight and still be rejected if the available certificate does not justify replacing the baseline.

The inverse heat experiment shows the same principle in a strongly ill-posed regime. All representative learned candidates are rejected at \(T=0.08\), even when the relative error is not large. This behaviour is consistent with the theory. As the final observation time increases, the forward heat equation smooths the initial condition more strongly, and the inverse problem becomes less stable. The certified radius reflects this loss of stability through the stability constant. A learned reconstruction may recover visible structure, but the no-harm rule requires the recovered structure to be supported by the certified radius before it can replace the baseline.

The stochastic residual experiment addresses a practical issue in physics-informed learning. Physics residuals are commonly estimated from collocation points, and a residual evaluated on the training set may be optimistic. The high-probability certificate uses independent validation collocation points to upper-bound the population residual. In the reported trials, the empirical mean squared residual fluctuates around the true mean squared residual, the high-probability upper bound decreases as the number of validation points increases, and the observed coverage rate is \(1.000\). This supports the use of independent residual validation as part of the no-harm workflow.

The main conceptual contribution is therefore not another physics-informed neural-network architecture. The contribution is a reusable decision rule for inverse learning: a learned solver should replace a baseline only when its residual-calibrated certificate dominates the baseline certificate. This principle is architecture-independent. It can wrap physics-informed neural networks, neural operators, unrolled inverse solvers, hybrid finite-element neural methods, classical PDE-constrained optimizers, or future inverse solvers that can report the residuals required for certification. In this sense, the framework changes the operational question from ``Did the learned reconstruction look good?'' to ``Was the learned reconstruction certified enough to replace the baseline?''

\subsection{Limitations}
\label{sec:limitations}

The framework depends on a stability estimate on the admissible class. This is not a weakness specific to the proposed method; it is a basic feature of inverse problems. Residuals cannot control reconstruction error unless the inverse map has enough stability on the class being considered. When the inverse problem is severely ill-posed, the stability constant may be large and the resulting radius may be conservative. This is visible in the inverse heat experiment, where the certificate becomes cautious as the final-time observation becomes less informative about the initial condition.

The certified radius also depends on the residual norms and their scaling. Data residuals, physics residuals, boundary or initial-condition residuals, and optimization residuals may have different units and magnitudes. The weights used in the operational certificate must therefore be fixed before evaluation and reported clearly. Poor scaling can make the certificate too permissive or too conservative. For this reason, the numerical validation reports the residual weights and treats the operational radius as a calibrated decision quantity.

The no-harm rule is not designed to maximize hidden ground-truth accuracy. It compares certified radii. A learned reconstruction may have lower true error than the baseline and still be rejected if its certified radius is larger. This is acceptable in a reliability framework because the true error is unavailable in real applications. The purpose is not to choose the reconstruction that would be best if the truth were known. The purpose is to prevent an uncertified learned output from silently replacing a safer certified baseline.

The stochastic residual result is stated for a fixed candidate reconstruction evaluated on an independent validation collocation set. A uniform certificate over a whole neural-network class would require additional capacity control, such as covering arguments, Rademacher complexity, PAC-Bayes bounds, or related uniform-convergence tools. The present result supports an a posteriori workflow: train the model first, then certify the computed reconstruction on independent residual-validation points.

The experiments are controlled numerical tests. They are designed to test the mechanism of the certificate across representative inverse-problem regimes rather than to exhaust all possible applications. Larger-scale validation in imaging, geophysics, and physical diagnostics should include realistic model error, heterogeneous noise, incomplete boundary information, uncertain observation operators, and operational constraints. Such studies would test how conservative the no-harm selector becomes in application-scale settings and how the certificate should be calibrated for different scientific domains.

\section{Conclusion}
\label{sec:conclusion}

This paper introduced a no-harm certification-and-selection framework for physics-informed inverse learning. The method attaches a residual-calibrated uncertainty radius to a learned inverse reconstruction and allows the learned output to replace a baseline only when its certified radius is no worse than the baseline radius, up to a prescribed tolerance. If this condition fails, the method falls back to the baseline rather than applying an uncertified correction or silently trusting the learned output.

The theoretical contribution is a residual-calibrated reliability principle. Under a conditional stability estimate on the admissible class, the data residual, physics residual, boundary or initial-condition residual, and noise level yield an a posteriori reconstruction-error bound. This gives a deterministic uncertainty radius. The no-harm theorem then turns this radius into a certificate-dominance rule: the selected reconstruction is controlled by the baseline certificate up to the chosen safety tolerance. The stochastic residual result extends the certificate to independently sampled validation collocation points, and the optimization residual records whether the learned solver has reached a meaningful stationarity level.

The numerical validation supports the principle. The sufficiency sweep identifies the operational threshold \(R_{\mathrm{learn}}/R_{\mathrm{base}}\le1\) in the strict no-harm setting. The Poisson and geophysical coefficient examples show that learned reconstructions are accepted when their certificates improve on the baseline. The shifted, hallucinated, and unfinished candidates show that unsupported learned outputs are rejected. The inverse heat example shows conservative behaviour in a strongly ill-posed regime. The limited-angle tomography example shows that the rule is based on available reliability evidence, not visual quality or hindsight error.

The main message is simple: physics-informed inverse learning should not only reconstruct. It should reconstruct, certify, and select. A learned reconstruction is sufficient to replace a baseline only when its residual-calibrated certificate supports that replacement. Otherwise, the safer output is the baseline.

\bibliographystyle{unsrt}
\bibliography{refs}

\appendix

\section{Proofs of theoretical results}
\label{app:theoretical-proofs}

This appendix gives the proofs of the theoretical results stated in the main text. The notation is the same as in Sections~\ref{sec:problem-setting}--\ref{sec:optimization}. Definitions and assumptions are not reproved. Each proof is labelled by the result it establishes.

\subsection{Proof of Theorem~\ref{thm:residual-error-bound}}

\begin{proof}[Proof of Theorem~\ref{thm:residual-error-bound}]
Assume \eqref{eq:observation} and Assumption~\ref{ass:conditional-stability}. Let \((\widehat u,\widehat q)\in\mathcal A\). Since
\[
y^\delta=\mathcal H u^\dagger+e,
\qquad
\|e\|_{\mathcal Y}\le \delta,
\]
we have
\[
\|y^\delta-\mathcal H u^\dagger\|_{\mathcal Y}
=
\|e\|_{\mathcal Y}
\le
\delta.
\]
By the triangle inequality in \(\mathcal Y\),
\begin{align}
\|\mathcal H\widehat u-\mathcal H u^\dagger\|_{\mathcal Y}
&\le
\|\mathcal H\widehat u-y^\delta\|_{\mathcal Y}
+
\|y^\delta-\mathcal H u^\dagger\|_{\mathcal Y} \notag\\
&\le
r_{\mathrm{data}}(\widehat u)+\delta .
\label{eq:app-data-bound}
\end{align}
Since \((\widehat u,\widehat q)\in\mathcal A\), Assumption~\ref{ass:conditional-stability} applies with \(u=\widehat u\) and \(q=\widehat q\). Hence
\begin{align}
\|\widehat q-q^\dagger\|_{\mathcal Q}
&\le
C_{\mathrm{stab}}
\left(
\|\mathcal H\widehat u-\mathcal H u^\dagger\|_{\mathcal Y}
+
\|\mathcal N(\widehat u,\widehat q)\|_{\mathcal Z}
+
\|\mathcal B(\widehat u,\widehat q)\|_{\mathcal B_0}
\right)^p .
\end{align}
Using \eqref{eq:app-data-bound}, together with
\[
\|\mathcal N(\widehat u,\widehat q)\|_{\mathcal Z}
=
r_{\mathrm{pde}}(\widehat u,\widehat q),
\qquad
\|\mathcal B(\widehat u,\widehat q)\|_{\mathcal B_0}
=
r_{\mathrm{bc}}(\widehat u,\widehat q),
\]
gives
\[
\|\widehat q-q^\dagger\|_{\mathcal Q}
\le
C_{\mathrm{stab}}
\left(
r_{\mathrm{data}}(\widehat u)
+
r_{\mathrm{pde}}(\widehat u,\widehat q)
+
r_{\mathrm{bc}}(\widehat u,\widehat q)
+
\delta
\right)^p ,
\]
which is \eqref{eq:aposteriori-error-bound}.
\end{proof}

\subsection{Proof of Corollary~\ref{cor:certified-containment}}

\begin{proof}[Proof of Corollary~\ref{cor:certified-containment}]
By Definition~\ref{def:uncertainty-radius},
\[
\mathcal U_\delta(\widehat u,\widehat q)
=
\left\{
\widetilde q\in\mathcal Q:
\|\widetilde q-\widehat q\|_{\mathcal Q}
\le
R_\delta(\widehat u,\widehat q)
\right\}.
\]
By Theorem~\ref{thm:residual-error-bound},
\[
\|\widehat q-q^\dagger\|_{\mathcal Q}
\le
R_\delta(\widehat u,\widehat q).
\]
Since norms are symmetric,
\[
\|q^\dagger-\widehat q\|_{\mathcal Q}
=
\|\widehat q-q^\dagger\|_{\mathcal Q}
\le
R_\delta(\widehat u,\widehat q).
\]
Thus \(q^\dagger\) satisfies the membership condition in \(\mathcal U_\delta(\widehat u,\widehat q)\), and therefore
\[
q^\dagger\in\mathcal U_\delta(\widehat u,\widehat q).
\]
\end{proof}

\subsection{Proof of Theorem~\ref{thm:no-harm-principle}}

\begin{proof}[Proof of Theorem~\ref{thm:no-harm-principle}]
Assume that \(A_{\mathrm{base}}(y^\delta)\) and \(A_{\mathrm{learn}}(y^\delta)\) are certificate-valid in the sense of Definition~\ref{def:certificate-valid-reconstruction}. Thus
\[
\|q_{\mathrm{base}}-q^\dagger\|_{\mathcal Q}
\le
\mathfrak R(A_{\mathrm{base}})
\]
and
\[
\|q_{\mathrm{learn}}-q^\dagger\|_{\mathcal Q}
\le
\mathfrak R(A_{\mathrm{learn}}).
\]
By Definition~\ref{def:no-harm-selector}, the selected output is
\[
(u_{\mathrm{safe}},q_{\mathrm{safe}})
=
S_{\mathrm{NH}}(y^\delta).
\]
There are two cases.

If the learned reconstruction is selected, then
\[
q_{\mathrm{safe}}=q_{\mathrm{learn}}.
\]
By certificate validity of \(A_{\mathrm{learn}}\),
\[
\|q_{\mathrm{safe}}-q^\dagger\|_{\mathcal Q}
=
\|q_{\mathrm{learn}}-q^\dagger\|_{\mathcal Q}
\le
\mathfrak R(A_{\mathrm{learn}}).
\]
Since the learned reconstruction is selected, the no-harm selector gives
\[
\mathfrak R(A_{\mathrm{learn}})
\le
\mathfrak R(A_{\mathrm{base}})
+
\varepsilon_{\mathrm{safe}}.
\]
Combining the last two inequalities gives
\[
\|q_{\mathrm{safe}}-q^\dagger\|_{\mathcal Q}
\le
\mathfrak R(A_{\mathrm{base}})
+
\varepsilon_{\mathrm{safe}}.
\]

If the learned reconstruction is rejected, then
\[
q_{\mathrm{safe}}=q_{\mathrm{base}}.
\]
By certificate validity of \(A_{\mathrm{base}}\),
\[
\|q_{\mathrm{safe}}-q^\dagger\|_{\mathcal Q}
=
\|q_{\mathrm{base}}-q^\dagger\|_{\mathcal Q}
\le
\mathfrak R(A_{\mathrm{base}}).
\]
This proves
\[
\|q_{\mathrm{safe}}-q^\dagger\|_{\mathcal Q}
\le
\mathfrak R(A_{\mathrm{base}}).
\]
Since
\[
\mathfrak R(A_{\mathrm{base}})
\le
\mathfrak R(A_{\mathrm{base}})
+
\varepsilon_{\mathrm{safe}},
\]
the bound
\[
\|q_{\mathrm{safe}}-q^\dagger\|_{\mathcal Q}
\le
\mathfrak R(A_{\mathrm{base}})
+
\varepsilon_{\mathrm{safe}}
\]
also holds in the rejected case. This proves both claims.
\end{proof}

\subsection{Proof of Theorem~\ref{thm:high-prob-residual}}

\begin{proof}[Proof of Theorem~\ref{thm:high-prob-residual}]
Fix an admissible pair \((u,q)\). Let \(X_1,\dots,X_M\) be independent samples from \(\nu\), and define
\[
Z_j
=
\left|\mathcal N(u,q)(X_j)\right|^2,
\qquad
j=1,\dots,M.
\]
By Assumption~\ref{ass:bounded-residual},
\[
0\le Z_j\le B_{\mathrm{pde}}
\qquad
\text{almost surely}.
\]
The empirical squared residual is
\[
\widehat r_{\mathrm{pde}}^{\,2}(u,q)
=
\frac1M\sum_{j=1}^M Z_j.
\]
The population squared residual is
\[
r_{\mathrm{pde},\nu}^{2}(u,q)
=
\int_\Omega
\left|\mathcal N(u,q)(x)\right|^2\,d\nu(x).
\]
Since each \(X_j\) has law \(\nu\),
\[
\mathbb E Z_j
=
\int_\Omega
\left|\mathcal N(u,q)(x)\right|^2\,d\nu(x)
=
r_{\mathrm{pde},\nu}^{2}(u,q).
\]
Thus
\[
\mathbb E\left[\widehat r_{\mathrm{pde}}^{\,2}(u,q)\right]
=
r_{\mathrm{pde},\nu}^{2}(u,q).
\]

By Hoeffding's inequality for independent bounded random variables \cite{Hoeffding1963}, for every \(t>0\),
\[
\mathbb P\left(
r_{\mathrm{pde},\nu}^{2}(u,q)
-
\widehat r_{\mathrm{pde}}^{\,2}(u,q)
\ge t
\right)
\le
\exp\left(
-\frac{2Mt^2}{B_{\mathrm{pde}}^2}
\right).
\]
Choose
\[
t
=
B_{\mathrm{pde}}
\sqrt{\frac{\log(1/\zeta)}{2M}},
\qquad
\zeta\in(0,1).
\]
Then
\[
\exp\left(
-\frac{2Mt^2}{B_{\mathrm{pde}}^2}
\right)
=
\exp\left(-\log(1/\zeta)\right)
=
\zeta.
\]
Therefore, with probability at least \(1-\zeta\),
\[
r_{\mathrm{pde},\nu}^{2}(u,q)
-
\widehat r_{\mathrm{pde}}^{\,2}(u,q)
\le
B_{\mathrm{pde}}
\sqrt{\frac{\log(1/\zeta)}{2M}}.
\]
Rearranging gives
\[
r_{\mathrm{pde},\nu}^{2}(u,q)
\le
\widehat r_{\mathrm{pde}}^{\,2}(u,q)
+
B_{\mathrm{pde}}
\sqrt{\frac{\log(1/\zeta)}{2M}},
\]
which is \eqref{eq:high-prob-residual}.
\end{proof}

\subsection{Proof of Corollary~\ref{cor:stochastic-containment}}

\begin{proof}[Proof of Corollary~\ref{cor:stochastic-containment}]
Assume that the conditional stability estimate in Assumption~\ref{ass:conditional-stability} is valid with the physics residual \(r_{\mathrm{pde},\nu}\). Let \((\widehat u,\widehat q)\in\mathcal A\) be fixed before drawing the independent validation collocation set.

Define the event
\[
\mathcal E_\zeta
=
\left\{
r_{\mathrm{pde},\nu}^{2}(\widehat u,\widehat q)
\le
\widehat r_{\mathrm{pde}}^{\,2}(\widehat u,\widehat q)
+
B_{\mathrm{pde}}
\sqrt{\frac{\log(1/\zeta)}{2M}}
\right\}.
\]
By Theorem~\ref{thm:high-prob-residual},
\[
\mathbb P(\mathcal E_\zeta)\ge 1-\zeta.
\]
On \(\mathcal E_\zeta\), the definition \eqref{eq:hp-pde-residual} gives
\[
r_{\mathrm{pde},\nu}(\widehat u,\widehat q)
\le
\widehat r_{\mathrm{pde},\zeta}(\widehat u,\widehat q).
\]
Applying Theorem~\ref{thm:residual-error-bound}, with \(r_{\mathrm{pde},\nu}\) in place of \(r_{\mathrm{pde}}\), gives
\[
\|\widehat q-q^\dagger\|_{\mathcal Q}
\le
C_{\mathrm{stab}}
\left(
r_{\mathrm{data}}(\widehat u)
+
r_{\mathrm{pde},\nu}(\widehat u,\widehat q)
+
r_{\mathrm{bc}}(\widehat u,\widehat q)
+
\delta
\right)^p .
\]
Since \(p\in(0,1]\) and all residual quantities are nonnegative, the map \(s\mapsto s^p\) is nondecreasing on \([0,\infty)\). Hence, on \(\mathcal E_\zeta\),
\[
\|\widehat q-q^\dagger\|_{\mathcal Q}
\le
C_{\mathrm{stab}}
\left(
r_{\mathrm{data}}(\widehat u)
+
\widehat r_{\mathrm{pde},\zeta}(\widehat u,\widehat q)
+
r_{\mathrm{bc}}(\widehat u,\widehat q)
+
\delta
\right)^p.
\]
By definition \eqref{eq:stochastic-radius}, the right-hand side is
\[
R_{\delta,\zeta}^{\mathrm{stoch}}(\widehat u,\widehat q).
\]
Therefore,
\[
\|\widehat q-q^\dagger\|_{\mathcal Q}
\le
R_{\delta,\zeta}^{\mathrm{stoch}}(\widehat u,\widehat q)
\]
with probability at least \(1-\zeta\).
\end{proof}

\subsection{Proof of Proposition~\ref{prop:optimization-contribution}}

\begin{proof}[Proof of Proposition~\ref{prop:optimization-contribution}]
By Assumption~\ref{ass:local-error-bound}, there exist a neighbourhood \(\mathcal V\) of \(\widehat\theta\), a stationary set \(\Theta^\star\), and a constant \(C_{\mathrm{opt}}>0\) such that
\[
\dist(\theta,\Theta^\star)
\le
C_{\mathrm{opt}}\|\nabla_\theta\mathcal L(\theta)\|
\qquad
\text{for all } \theta\in\mathcal V.
\]
Since \(\widehat\theta\in\mathcal V\), evaluating this inequality at \(\theta=\widehat\theta\) gives
\[
\dist(\widehat\theta,\Theta^\star)
\le
C_{\mathrm{opt}}\|\nabla_\theta\mathcal L(\widehat\theta)\|.
\]
By definition \eqref{eq:opt-residual-repeat},
\[
r_{\mathrm{opt}}(\widehat\theta)
=
\|\nabla_\theta\mathcal L(\widehat\theta)\|.
\]
Therefore
\[
\dist(\widehat\theta,\Theta^\star)
\le
C_{\mathrm{opt}}r_{\mathrm{opt}}(\widehat\theta),
\]
which proves \eqref{eq:opt-distance-bound}.

Moreover, if \(\alpha_{\mathrm{opt}}>0\), then
\[
\alpha_{\mathrm{opt}}r_{\mathrm{opt}}(\widehat\theta)
\ge
\frac{\alpha_{\mathrm{opt}}}{C_{\mathrm{opt}}}
\dist(\widehat\theta,\Theta^\star).
\]
Thus the optimization term in the total certificate \eqref{eq:total-certificate} lower-bounds the distance-to-stationarity up to the constant \(\alpha_{\mathrm{opt}}/C_{\mathrm{opt}}\). If \(r_{\mathrm{opt}}(\widehat\theta)\) is large, then the certificate records that the computed reconstruction has not been certified as an accurately optimized physics-informed solution. This establishes the stated contribution of the optimization residual to the certificate.
\end{proof}


\section{Full validation tables}
\label{app:full-validation-results}

This appendix reports the numerical outputs used in Section~\ref{sec:numerical-validation}. The main text reports compact decision summaries in Table~\ref{tab:main-noharm-summary}. Table~\ref{tab:appendix-noharm-full} reports the candidate-level no-harm decisions. Table~\ref{tab:appendix-stability-constants} reports the stability and conditioning quantities used to compute the certified radii. Table~\ref{tab:appendix-stochastic-sweep} reports the stochastic residual sweep. The symbol ``--'' indicates that the corresponding quantity is not applicable to that experiment, not that the value is missing.

\subsection{Candidate-level no-harm decisions}

\begin{table}[!t]
\centering
\scriptsize
\caption{Candidate-level no-harm decisions across all validation experiments. \(R_{\mathrm{base}}\) is the certified radius of the baseline and \(R_{\mathrm{learn}}\) is the certified radius of the learned candidate. The decision is computed using \(R_{\mathrm{learn}}\le R_{\mathrm{base}}+\varepsilon_{\mathrm{safe}}\).}
\label{tab:appendix-noharm-full}
\resizebox{\textwidth}{!}{%
\begin{tabular}{llcrrrll}
	\toprule
	\textbf{Experiment} & \textbf{Scenario} & \textbf{Candidate} & \(\boldsymbol{R_{\mathrm{base}}}\) & \(\boldsymbol{R_{\mathrm{learn}}}\) & \textbf{Rel. error} & \textbf{Decision} & \textbf{Safe output} \\
	\midrule
	Poisson source & default & learned good & 5.375 & 4.930 & 0.0719 & accept & learned \\
	Poisson source & default & shifted learned & 5.375 & 19.28 & 0.5739 & reject & baseline \\
	Poisson source & default & unfinished PINN & 5.375 & 9045 & 0.1020 & reject & baseline \\
	Inverse heat & \(T=0.02\) & learned good & 0.0312 & 0.1486 & 0.1059 & reject & baseline \\
	Inverse heat & \(T=0.02\) & hallucinated high freq. & 0.0312 & 0.6688 & 0.9633 & reject & baseline \\
	Inverse heat & \(T=0.02\) & shifted learned & 0.0312 & 0.9933 & 0.7599 & reject & baseline \\
	Inverse heat & \(T=0.08\) & learned good & 0.0403 & 0.1181 & 0.0748 & reject & baseline \\
	Inverse heat & \(T=0.08\) & hallucinated high freq. & 0.0403 & 0.1612 & 0.9604 & reject & baseline \\
	Inverse heat & \(T=0.08\) & shifted learned & 0.0403 & 1.116 & 0.7599 & reject & baseline \\
	Inverse heat & \(T=0.16\) & learned good & 0.0464 & 0.2388 & 0.1281 & reject & baseline \\
	Inverse heat & \(T=0.16\) & hallucinated high freq. & 0.0464 & 0.2389 & 0.9660 & reject & baseline \\
	Inverse heat & \(T=0.16\) & shifted learned & 0.0464 & 1.307 & 0.7599 & reject & baseline \\
	Tomography & default & learned good & 6.520 & 7.922 & 0.0190 & reject & baseline \\
	Tomography & default & hallucinated learned & 6.520 & 76.93 & 0.3204 & reject & baseline \\
	Geophysical coeff. & default & learned good & 6.048 & 4.524 & 0.0050 & accept & learned \\
	Geophysical coeff. & default & shifted learned & 6.048 & 11.01 & 0.0635 & reject & baseline \\
	\bottomrule
\end{tabular}%
}
\end{table}

\subsection{Stability and conditioning quantities}

\begin{table}[!ht]
\centering
\caption{Stability and conditioning quantities used in the certified-radius calculations. The symbol ``--'' means that the quantity is not part of that experiment. For example, \(n_{\mathrm{angles}}\) applies only to tomography, while \(T\) and \(\kappa\) apply only to the inverse heat experiment.}
\label{tab:appendix-stability-constants}
\resizebox{\linewidth}{!}{
\begin{tabular}{llrrrrrrrrrr}
	\toprule
	\textbf{Experiment} & \textbf{Scenario} & \(\boldsymbol{\sigma_{\min}}\) & \(\boldsymbol{C_{\mathrm{stab}}}\) & \(\boldsymbol{n_{\mathrm{basis}}}\) & \(\boldsymbol{n_{\mathrm{obs}}}\) & \(\boldsymbol{T}\) & \(\boldsymbol{\kappa}\) & \(\boldsymbol{n_{\mathrm{meas}}}\) & \(\boldsymbol{n_{\mathrm{angles}}}\) & \(\boldsymbol{\kappa(F)}\) & \(\boldsymbol{\kappa(J)}\) \\
	\midrule
	Poisson source & default & 0.0005441 & 1838 & 10 & 35 & -- & -- & -- & -- & 99.55 & -- \\
	Inverse heat & \(T=0.02\) & 0.5719 & 1.748 & 8 & 45 & 0.02 & 0.004 & -- & -- & 1.063 & -- \\
	Inverse heat & \(T=0.08\) & 0.4923 & 2.031 & 8 & 45 & 0.08 & 0.004 & -- & -- & 1.231 & -- \\
	Inverse heat & \(T=0.16\) & 0.4026 & 2.484 & 8 & 45 & 0.16 & 0.004 & -- & -- & 1.501 & -- \\
	Tomography & default & 0.4653 & 2.149 & 36 & -- & -- & -- & 420 & 15 & 40.92 & -- \\
	Geophysical coeff. & default & 0.002414 & 414.2 & 6 & 30 & -- & -- & -- & -- & -- & 23.90 \\
	\bottomrule
\end{tabular}}
\end{table}

\subsection{Stochastic residual sweep}

\begin{table}[!ht]
\centering
\caption{Full stochastic residual sweep. The table reports the empirical mean squared residual and the high-probability upper bound over 250 repetitions at confidence parameter \(\zeta=0.05\).}
\label{tab:appendix-stochastic-sweep}
\begin{tabular}{rrrrrrr}
	\toprule
	\(\boldsymbol{M}\) & \textbf{Reps.} & \textbf{True MSR} & \textbf{Empirical MSR} & \textbf{HP upper bound} & \textbf{Coverage} & \(\boldsymbol{\zeta}\) \\
	\midrule
	20 & 250 & 0.2252 & 0.2313 & 1.075 & 1.000 & 0.05 \\
	50 & 250 & 0.2252 & 0.2279 & 0.7618 & 1.000 & 0.05 \\
	100 & 250 & 0.2252 & 0.2209 & 0.5984 & 1.000 & 0.05 \\
	200 & 250 & 0.2252 & 0.2245 & 0.4914 & 1.000 & 0.05 \\
	500 & 250 & 0.2252 & 0.2238 & 0.3926 & 1.000 & 0.05 \\
	1000 & 250 & 0.2252 & 0.2243 & 0.3437 & 1.000 & 0.05 \\
	2000 & 250 & 0.2252 & 0.2237 & 0.3081 & 1.000 & 0.05 \\
	\bottomrule
\end{tabular}
\end{table}

\section{Implementation details}
\label{app:implementation-details}

All experiments were run with fixed random seeds. The validation scripts generate the synthetic inverse-problem instances, compute baseline and learned candidates, evaluate the residual components, compute certified radii, apply the no-harm rule, and export the tables and figures used in Section~\ref{sec:numerical-validation}. The code records the noise level, stability constant, residual components, candidate errors, certified radii, coverage indicators, and no-harm decisions for each experiment.

The figures in the main text are produced from the exported reconstruction and diagnostic files. The tables in Appendix~\ref{app:full-validation-results} are produced directly from the exported CSV files. No table values are manually altered. Numerical values are rounded for presentation.

\section{Notes on certified radii}
\label{app:certified-radius-notes}

The certified radii reported in the numerical experiments are computed from the residual-calibrated formula developed in Section~\ref{sec:conditional-stability}. The constants are problem-dependent and are reported explicitly in Table~\ref{tab:appendix-stability-constants}. These constants should be interpreted as stability and conditioning quantities used for transparent numerical certification, not as sharp optimal constants for each inverse problem.

A smaller true reconstruction error does not necessarily imply a smaller certified radius. The radius measures reliability under the available data, residuals, noise level, and stability estimate. This is why some learned candidates with low true error may still be rejected. Such rejection is conservative, but it is consistent with the no-harm principle.

\end{document}
